\documentclass[11pt]{article}
\usepackage[margin=1in]{geometry}
\usepackage{amsmath,amssymb,amsthm}
\usepackage[T1]{fontenc}
\usepackage{lmodern}
\usepackage{xcolor}
\usepackage{listings}
\usepackage{booktabs}
\usepackage{longtable}
\usepackage{array}
\usepackage[hidelinks]{hyperref}
\usepackage{url}

\newtheorem{theorem}{Theorem}
\newtheorem{lemma}[theorem]{Lemma}
\newtheorem{corollary}[theorem]{Corollary}
\newtheorem{proposition}[theorem]{Proposition}
\theoremstyle{definition}
\newtheorem{definition}[theorem]{Definition}
\newtheorem{example}[theorem]{Example}
\theoremstyle{remark}
\newtheorem{remark}[theorem]{Remark}

% Code-like text (not checked by a name gate).
\newcommand{\lean}[1]{\texttt{#1}}
% A declaration of one of the five developments. Every \decl{...} is meant to be checked by a
% name gate (grep for a declaration of that name in the five source trees).
\DeclareUrlCommand\decl{\urlstyle{tt}\def\UrlBreaks{\do\_\do\.}}
% A Mathlib declaration (pinned revision 0df444a3); not one of ours.
\DeclareUrlCommand\mdecl{\urlstyle{tt}\def\UrlBreaks{\do\_\do\.}}
\DeclareUnicodeCharacter{2081}{\ensuremath{_1}}
\DeclareUnicodeCharacter{2082}{\ensuremath{_2}}
\DeclareUnicodeCharacter{2083}{\ensuremath{_3}}
\DeclareUnicodeCharacter{211D}{\ensuremath{\mathbb{R}}}
\DeclareUnicodeCharacter{03C9}{\ensuremath{\omega}}
\newcommand{\F}{\mathbb{F}}
\newcommand{\RR}{\mathbb{R}}
\newcommand{\CC}{\mathbb{C}}
\newcommand{\QQ}{\mathbb{Q}}
\newcommand{\PP}{\mathbb{P}}
% Outputs that a separate job is producing. Must all be filled in before submission.
\newcommand{\PH}[1]{\textcolor{red}{\textbf{[PLACEHOLDER:} #1\textbf{]}}}

% version-2 commits: desargues=086ca718192df1cb919b72a7e79d9fb26e43e723 monge=82dd27d7b81e60e2aa57873e98cc6fa0dce0728f simson=e3d2c46abd6a6da2a4571a28cebc47c8447a0618 napoleon=0829a1483687956644b6b5eb3eb26205c6d7e9b7 lami=e692d1a9bc0004444e33b76132275461a2e26cf2
% Commit hashes of the published repositories (content commits; checked against the build machine by SHA-256).
\newcommand{\commitdesargues}{\href{\repo{desargues}/tree/b6409255c844abb6e1cad80b23f8182fc3934116}{\texttt{b640925}}}
\newcommand{\commitmonge}{\href{\repo{monge}/tree/a01f7504d12fd492dd30c66f880905b21cca3895}{\texttt{a01f750}}}
\newcommand{\commitsimson}{\href{\repo{simson}/tree/7c239b31600575db241b45dda4930cf9486e7d17}{\texttt{7c239b3}}}
\newcommand{\commitnapoleon}{\href{\repo{napoleon}/tree/e16592ae04be8e5516a2caef52418c6c60df5165}{\texttt{e16592a}}}
\newcommand{\commitlami}{\href{\repo{lami}/tree/ea6b038b1a3eebd9b02465650901f4b5d4ddaf19}{\texttt{ea6b038}}}
% Plain-text forms of the same, used inside URLs (replace together with the above).
\newcommand{\urlcommitdesargues}{b6409255c844abb6e1cad80b23f8182fc3934116}
\newcommand{\urlcommitmonge}{a01f7504d12fd492dd30c66f880905b21cca3895}
\newcommand{\urlcommitsimson}{7c239b31600575db241b45dda4930cf9486e7d17}
\newcommand{\urlcommitnapoleon}{e16592ae04be8e5516a2caef52418c6c60df5165}
\newcommand{\urlcommitlami}{ea6b038b1a3eebd9b02465650901f4b5d4ddaf19}
\expandafter\def\csname commitof:desargues\endcsname{\urlcommitdesargues}
\expandafter\def\csname commitof:monge\endcsname{\urlcommitmonge}
\expandafter\def\csname commitof:simson\endcsname{\urlcommitsimson}
\expandafter\def\csname commitof:napoleon\endcsname{\urlcommitnapoleon}
\expandafter\def\csname commitof:lami\endcsname{\urlcommitlami}
\newcommand{\commitof}[1]{\csname commitof:#1\endcsname}
\newcommand{\repo}[1]{https://github.com/jpbald93/#1}
% The repositories are published (5 Oct 2026); \src links to the cited commits. With
% \publishedfalse, \src prints plain text instead.
\newif\ifpublished
\publishedtrue
% \src{repo}{Dir}{File}{n}: line n of Dir/File.lean in the given repository.
\ifpublished
\newcommand{\src}[4]{\href{\repo{#1}/blob/\commitof{#1}/#2/#3.lean\#L#4}{\texttt{#3}:#4}}
\else
\newcommand{\src}[4]{\texttt{#3}:#4}
\fi

\lstdefinelanguage{lean}{
  morekeywords={theorem,lemma,def,class,example,by,fun,have,show,exact,intro,namespace,end,import,private,variable,noncomputable},
  sensitive=true, morecomment=[l]{--}, morecomment=[s]{/-}{-/}, morestring=[b]"}
\lstset{language=lean, basicstyle=\ttfamily\small, keywordstyle=\bfseries,
  columns=fullflexible, keepspaces=true, breaklines=true, frame=single, framerule=0.3pt,
  literate={≠}{{$\neq$}}1 {↔}{{$\leftrightarrow$}}2 {→}{{$\to$}}1 {↑}{{$\uparrow$}}1
           {⨯₃}{{$\times_3$}}2 {⬝ᵥ}{{$\cdot_v$}}2 {₁}{{$_1$}}1 {₂}{{$_2$}}1 {₃}{{$_3$}}1
           {₄}{{$_4$}}1 {ℓ}{{$\ell$}}1 {≤}{{$\le$}}1 {ℙ}{{$\mathbb{P}$}}1 {ℂ}{{$\mathbb{C}$}}1
           {ℝ}{{$\mathbb{R}$}}1 {ℚ}{{$\mathbb{Q}$}}1 {ω}{{$\omega$}}1 {‖}{{$\|$}}1
           {∀}{{$\forall$}}1 {∃}{{$\exists$}}1 {∧}{{$\wedge$}}1 {∨}{{$\vee$}}1
           {¬}{{$\neg$}}1 {∈}{{$\in$}}1 {∡}{{$\measuredangle$}}1 {∠}{{$\angle$}}1
           {π}{{$\pi$}}1 {γ}{{$\gamma$}}1 {'}{{$^{\prime}$}}1 {×}{{$\times$}}1}

\title{Classical plane geometry by coordinate certificates:\\
Desargues, Monge, Simson, Napoleon and Lami in Lean~4}
\author{Josh Bald\\ \small Independent researcher, Ontario, Canada\\ \small\texttt{jpbald93@gmail.com}\\
\small ORCID: \texttt{0009-0002-1317-6489}
\thanks{\copyright~2026 Josh Bald. Licensed under CC BY 4.0 (\url{https://creativecommons.org/licenses/by/4.0/}). The accompanying Lean and Python code is licensed under the Apache License 2.0.}}
\date{7 October 2026; version~2, 22 October 2026}

\emergencystretch=3em
\begin{document}
\maketitle

\begin{abstract}
We formalise five classical theorems of plane geometry in Lean~4 with Mathlib: Desargues's
theorem in the projective plane, together with its converse; Monge's theorem that the three
external centres of similitude of three circles are collinear; the Wallace--Simson theorem on the
pedal triangle of a point of the circumcircle; Napoleon's theorem for the outer and the inner
Napoleon triangles; and Lami's theorem for three forces in equilibrium. Every proof has the same
core. The theorem is reduced to an explicit polynomial or rational identity in coordinates, which
we call a \emph{certificate}, and the identity is closed by \lean{ring}, \lean{field\_simp} or
\lean{linear\_combination} with a stated multiplier. Around this core, each development makes
precise which non-degeneracy conditions are needed, which are not, and why. Desargues's theorem and
Monge's theorem are proved over an arbitrary field, and the Desargues identity over an arbitrary
commutative ring. The Simson line and Napoleon's theorem are proved in the Euclidean plane,
modelled as~$\CC$. Lami's theorem is proved through the oriented area form of a two-dimensional
oriented real inner product space, and in unoriented form in every real inner product space. The
five libraries, without the version-2 files described next, have about 1{,}430 lines of Lean, use no \lean{sorry}, and depend only on Lean's
three standard axioms. Version~2 adds to each library a file giving Mathlib-only formulations of the main theorems
(for Napoleon's theorem, of its non-degenerate configurations), each with a non-degenerate witness, and a file proving from the libraries
textbook statements written in Lean by an AI model that saw neither our code nor this paper; these
include the geometric converse of the Wallace--Simson theorem, which the first version did not prove. The mathematics is classical and the method is old, and several of the
results are not new to formal mathematics (Section~\ref{sec:related}). Monge's three-circle
theorem has two earlier, independent and complete Lean proofs, one of which predates ours by four
months; Napoleon's theorem has
earlier formalisations in Lean, Isabelle and Rocq, and a later independent one in Lean; the Wallace--Simson theorem and its converse
are formalised in Isabelle and in Rocq; Desargues's theorem is formalised in at least five other
systems, including a coordinate (bracket) proof in HOL~Light; and an open, unmerged Mathlib pull
request proves the affine (parallel) form of Desargues's theorem. We found no earlier Lean proof
of the projective-plane form of Desargues's theorem, of the Wallace--Simson theorem, or of Lami's
theorem; these searches are not exhaustive.
\end{abstract}

\section{Introduction}

Elementary plane geometry is a natural testing ground for proof assistants: the statements are
short and familiar, but a faithful formal statement has to say exactly which configurations are
allowed. Textbooks rarely do. ``Two triangles in perspective from a point'' may or may not
allow coincident vertices; ``three circles'' may or may not allow equal radii; ``the ratio
$|F_1|/\sin\alpha$'' is not defined when $\sin\alpha=0$. A formal statement must decide each of
these questions, and the decision is part of the result.

The proofs in this paper are coordinate proofs in the tradition of mechanical geometry theorem
proving~\cite{Chou}. Each theorem is reduced to an identity of the form
\[
  (\text{conclusion})\cdot(\text{non-degenerate factor}) \;=\; \sum_i c_i\cdot(\text{hypothesis}_i),
\]
with the multipliers $c_i$ written out (or with no hypotheses at all on the right), and the
identity is checked by a decision procedure for commutative rings. We call such an identity a
\emph{certificate}. The certificate is only the algebraic core. Each development also has a
translation layer connecting the coordinate statement with a standard notion of Mathlib, such as
\mdecl{Projectivization.IsCollinear}, \mdecl{Collinear}, \mdecl{dist} or
\mdecl{Orientation.oangle}. That layer is where most of the non-degeneracy hypotheses enter. The
point of the paper is the analysis of those hypotheses as much as the identities.

\paragraph{Contributions.}
\begin{itemize}
\item Desargues's theorem in Mathlib's projective plane $\PP(K^3)$ over an arbitrary field $K$,
in a determinant form that needs only a nonzero centre of perspective, and in a projective form
whose distinctness hypotheses are exactly those that make Mathlib's \mdecl{cross} compute sides
and intersection points (Section~\ref{sec:desargues}). The converse is proved in both forms. The
underlying factorisation $[PQR]=[abc]\,[a'b'c']\,[AA'\,BB'\,CC']$ is proved over every
commutative ring.
\item Monge's theorem in coordinates over an arbitrary field, with pairwise distinct radii as the
only hypothesis, and no hypothesis on the centres (Section~\ref{sec:monge}). This theorem is
already proved in Lean, twice and independently; our version differs only in being stated over an
arbitrary field and in its explicit analysis of the hypotheses.
\item The forward direction of the Wallace--Simson theorem in $\CC$, with conclusion Mathlib's
\mdecl{Collinear} and hypotheses ``the vertices are pairwise distinct'' and Mathlib's
\mdecl{EuclideanGeometry.Cospherical}, together with a determinant-level equivalence for a
non-degenerate triangle. The first version did not prove the geometric converse; version~2 proves
it, in \lean{EuclideanSpace ℝ (Fin 2)}, in a Layer-15 file (\decl{Simson.Classical.cov_A3_b};
Sections~\ref{sec:simson} and~\ref{sec:l1415}). This
theorem, with its converse, is already formalised in Isabelle and in Rocq.
\item Napoleon's theorem in $\CC$ as an equality of distances, for both Napoleon triangles and
with no non-degeneracy hypothesis, together with a signed-area lemma identifying which triangles
are erected outward (Section~\ref{sec:napoleon}).
\item Lami's theorem in oriented and unoriented forms, product and ratio forms, an explicit
treatment of the degenerate (parallel) case, and the oriented converse
(Section~\ref{sec:lami}).
\item A table of the non-degeneracy conditions of all five theorems
(Table~\ref{tab:nondeg}), recording for each whether the certificate needs it, the translation
needs it, or it is only needed for the statement to mean what it says.
\item (Version~2.) Mathlib-only formulations of the six numbered theorems (for Napoleon, of the
non-degenerate configurations), each with a non-degenerate witness, and proofs from the libraries of textbook statements of all five theorems
written in Lean by a model that saw neither code nor paper (Section~\ref{sec:l1415}).
\end{itemize}

\paragraph{Availability.}
The five developments are publicly available at the following addresses:
\path{github.com/jpbald93/desargues},
\path{github.com/jpbald93/monge},
\path{github.com/jpbald93/simson},
\path{github.com/jpbald93/napoleon} and
\path{github.com/jpbald93/lami}.
Line references (Table~\ref{tab:thms}) are to commits \commitdesargues, \commitmonge,
\commitsimson, \commitnapoleon\ and \commitlami\ respectively, and link to those commits. In the
first four repositories that commit adds a module \path{Bridge.lean} of definition bridges
(Section~\ref{sec:correct}) on top of the original content commit and a citation-only commit;
apart from one added \lean{import} line at the end of each root file, the earlier files are
unchanged. The Lami commit changes only the citation in \path{README.md}, so its Lean sources are
those of its content commit \texttt{a7a039b}. Each repository has since received three further
commits: one correcting the citation in \path{README.md}, one making the verification scripts
\path{gate.sh} and \path{tests/tamper.sh} portable to BSD-derived systems, which had required GNU
\path{mktemp}, \path{sed} and \path{python3}, and one adding licence and citation metadata
(\path{LICENSE}, \path{CITATION.cff}, \path{README.md}); none changes any Lean source, so the commits
cited above remain the sources of the formalisations. Version~2 of this paper adds the Layer-14 and
Layer-15 files (Section~\ref{sec:l1415}) in one new commit of each repository (desargues
\texttt{086ca718192df1cb919b72a7e79d9fb26e43e723}, monge \texttt{82dd27d7b81e60e2aa57873e98cc6fa0dce0728f}, simson \texttt{e3d2c46abd6a6da2a4571a28cebc47c8447a0618}, napoleon \texttt{0829a1483687956644b6b5eb3eb26205c6d7e9b7},
lami \texttt{e692d1a9bc0004444e33b76132275461a2e26cf2}), published on GitHub. That commit adds \path{L14.lean} and \path{L15Classical.lean} to the
library directory (for example \path{Simson/L15Classical.lean}), appends their imports to the root
file, adds the new declarations to the list checked by \path{gate.sh}, updates \path{README.md}, and
removes the \path{scratch/} directory, which is not part of the library. The other library files are
unchanged, so the statements, files and line references of the commits cited above are unchanged. The code in all five repositories (Lean
sources, gate and tamper scripts, and the Python certificate scripts) is released under the Apache
License 2.0, the licence used by Mathlib; this paper is released under CC BY 4.0. The
toolchain is Lean \texttt{v4.33.1}~\cite{Lean4}
with Mathlib~\cite{Mathlib} revision \texttt{0df444a3} (tag \texttt{v4.33.1}), pinned in each
\path{lake-manifest.json}. The libraries and their sizes are \lean{Desargues} (five files,
452 lines), \lean{Monge} (two files, 173 lines), \lean{Simson} (four files, 320 lines),
\lean{Napoleon} (two files, 205 lines) and \lean{Lami} (one file, 284 lines), each with a root file
that only imports. Version~2 adds 218 and 258 lines (Desargues), 91 and 261 (Monge), 121 and 168
(Simson), 334 and 503 (Napoleon), and 73 and 98 (Lami) in \path{L14.lean} and
\path{L15Classical.lean}. All declarations are in the namespace of the library name (the Layer-15
ones in its sub-namespace, e.g.\ \lean{Simson.Classical}). Table~\ref{tab:thms}
lists the main formal statements.
Supporting material (the Layer-14 and Layer-15 reports, the blind writer's statements and report, the transcript audit, and the reviews of this version with the fixes made in response) is in the folder \texttt{PaperA-coordinate-certificates} of the repository \url{https://github.com/jpbald93/LEAN-verification}.

{\small
\setlength{\LTcapwidth}{\textwidth}
\begin{longtable}{@{}>{\raggedright\arraybackslash}p{0.28\textwidth}
  >{\raggedright\arraybackslash}p{0.31\textwidth}
  >{\raggedright\arraybackslash}p{0.13\textwidth}
  >{\raggedright\arraybackslash}p{0.19\textwidth}@{}}
\caption{Informal statements and their formal counterparts. Line numbers refer to the first line
of the declaration. ``Field'' means an arbitrary field~$K$.}\label{tab:thms}\\
\toprule
Informal statement & Lean name & Source & Setting / hypotheses \\ \midrule
\endfirsthead
\toprule
Informal statement & Lean name & Source & Setting / hypotheses \\ \midrule
\endhead
\bottomrule
\endfoot
\multicolumn{4}{@{}l}{\emph{Desargues (Section~\ref{sec:desargues})}}\\
Bracket $[uvw]=u\cdot(v\times w)$ & \decl{Desargues.br} & \src{desargues}{Desargues}{Basic}{26} & comm.\ ring; definition \\
$[uvw]=\det(u,v,w)$ & \decl{Desargues.br_eq_det} & \src{desargues}{Desargues}{Basic}{28} & comm.\ ring; none \\
Certificate $[PQR]=[abc][a'b'c'][AA'BB'CC']$ & \decl{Desargues.desargues_det_identity} & \src{desargues}{Desargues}{Basic}{34} & comm.\ ring; none \\
Bracket $0$ $\Rightarrow$ collinear in $\PP^2$ & \decl{Desargues.isCollinear_mk_of_br_eq_zero} & \src{desargues}{Desargues}{Plane}{27} & field; $u,v,w\ne0$, $[uvw]=0$ \\
Collinear in $\PP^2$ $\Rightarrow$ bracket $0$ & \decl{Desargues.br_eq_zero_of_isCollinear_mk} & \src{desargues}{Desargues}{Plane}{59} & field; $u,v,w\ne0$, \mdecl{IsCollinear} \\
Common orthogonal $o\ne0$ $\Rightarrow$ bracket $0$ & \decl{Desargues.br_eq_zero_of_dotProduct_eq_zero} & \src{desargues}{Desargues}{Plane}{88} & field; $o\ne0$, $o\cdot x=o\cdot y=o\cdot z=0$ \\
Bracket $0$ $\Rightarrow$ common orthogonal $o\ne 0$ & \decl{Desargues.exists_dotProduct_eq_zero_of_br_eq_zero} & \src{desargues}{Desargues}{Plane}{100} & field; $[xyz]=0$ \\
Desargues, determinant form (Thm.~\ref{thm:desdet}) & \decl{Desargues.desargues_det} & \src{desargues}{Desargues}{Main}{33} & field; $o\ne0$, $[oaa']=[obb']=[occ']=0$ \\
Desargues in $\PP^2$ (Thm.~\ref{thm:desmk}) & \decl{Desargues.desargues_mk} & \src{desargues}{Desargues}{Main}{59} & field, decidable eq.; $a,b,c,a',b',c',o\ne0$; six vertex inequalities $A\ne B$, $B\ne C$, $C\ne A$, $A'\ne B'$, $B'\ne C'$, $C'\ne A'$; $[oaa']=[obb']=[occ']=0$; $AB\ne A'B'$, $BC\ne B'C'$, $CA\ne C'A'$ \\
Converse, determinant form & \decl{Desargues.desargues_converse_det} & \src{desargues}{Desargues}{Main}{88} & field; $[abc]\ne0$, $[a'b'c']\ne0$, $[PQR]=0$ \\
Converse in $\PP^2$ (dual form) & \decl{Desargues.desargues_converse_mk} & \src{desargues}{Desargues}{Main}{106} & field, decidable eq.; $a\times a',b\times b',c\times c'\ne0$; $[abc],[a'b'c']\ne0$; $P,Q,R\ne0$; $\{P,Q,R\}$ \mdecl{IsCollinear} \\
Non-degeneracy of the example & \decl{Desargues.ex_nonvacuous} & \src{desargues}{Desargues}{Examples}{41} & $K=\QQ$; none \\
Desargues applied to the example & \decl{Desargues.ex_desargues} & \src{desargues}{Desargues}{Examples}{74} & $K=\QQ$; none \\
\multicolumn{4}{@{}l}{\emph{Monge (Section~\ref{sec:monge})}}\\
External centre of similitude & \decl{Monge.ext} & \src{monge}{Monge}{Basic}{31} & field; definition \\
Signed double area & \decl{Monge.det3} & \src{monge}{Monge}{Basic}{36} & field; definition \\
Monge's theorem (Thm.~\ref{thm:monge}) & \decl{Monge.monge} & \src{monge}{Monge}{Basic}{42} & field; $r_1\ne r_2$, $r_2\ne r_3$, $r_3\ne r_1$ \\
Non-vacuity example over $\QQ$ & \decl{Monge.monge_example} & \src{monge}{Monge}{Basic}{54} & $K=\QQ$; none \\
Values of the example's centres & \decl{Monge.monge_example_values} & \src{monge}{Monge}{Basic}{61} & $K=\QQ$; none \\
\multicolumn{4}{@{}l}{\emph{Simson (Section~\ref{sec:simson})}}\\
Foot of perpendicular & \decl{Simson.perpFoot} & \src{simson}{Simson}{Basic}{28} & $\CC$; definition \\
Foot lies on the line & \decl{Simson.perpFoot_mem_line} & \src{simson}{Simson}{Basic}{50} & $\CC$; none \\
Foot is orthogonal & \decl{Simson.inner_perpFoot} & \src{simson}{Simson}{Basic}{54} & $\CC$; $u\ne v$ \\
Concyclicity determinant & \decl{Simson.concyc} & \src{simson}{Simson}{Identity}{16} & $\CC$; definition \\
Concyclic $\Rightarrow$ determinant $0$ & \decl{Simson.concyc_eq_zero_of_dist} & \src{simson}{Simson}{Identity}{27} & $\CC$; $|a-o|=|b-o|=|c-o|=|p-o|=r$ \\
Certificate (pedal area identity) & \decl{Simson.area2_perpFoot_mul} & \src{simson}{Simson}{Identity}{43} & $\CC$; $a\ne b$, $b\ne c$, $c\ne a$ \\
Simson, area form & \decl{Simson.simson_area} & \src{simson}{Simson}{Main}{51} & $\CC$; $a,b,c$ pairwise distinct; $|a-o|=|b-o|=|c-o|=|p-o|=r$ \\
Simson line, forward direction (Thm.~\ref{thm:simson}) & \decl{Simson.simson} & \src{simson}{Simson}{Main}{64} & $\CC$; $a,b,c$ pairwise distinct; $\{a,b,c,p\}$ \mdecl{Cospherical} \\
Determinant-level equivalence (not the geometric converse) & \decl{Simson.area2_perpFoot_eq_zero_iff} & \src{simson}{Simson}{Main}{74} & $\CC$; $\mathrm{area2}(a,b,c)\ne0$ \\
\multicolumn{4}{@{}l}{\emph{Napoleon (Section~\ref{sec:napoleon})}}\\
Apex of erected triangle & \decl{Napoleon.apex} & \src{napoleon}{Napoleon}{Basic}{28} & $\CC$; definition \\
Erected triangle is equilateral & \decl{Napoleon.apex_equilateral} & \src{napoleon}{Napoleon}{Basic}{64} & $|\omega|=1$, $|1-\omega|=1$ \\
Certificate $G_b-G_a=\omega(G_c-G_a)$ & \decl{Napoleon.key} & \src{napoleon}{Napoleon}{Basic}{74} & $\omega^2-\omega+1=0$ \\
Napoleon, general $\omega$ & \decl{Napoleon.napoleon_general} & \src{napoleon}{Napoleon}{Basic}{81} & $\omega^2-\omega+1=0$, $|\omega|=1$, $|1-\omega|=1$ (redundant, see text); no non-degeneracy \\
Outer Napoleon triangle (Thm.~\ref{thm:napoleon}) & \decl{Napoleon.napoleon_outer} & \src{napoleon}{Napoleon}{Basic}{98} & none \\
Inner Napoleon triangle & \decl{Napoleon.napoleon_inner} & \src{napoleon}{Napoleon}{Basic}{104} & none \\
Signed area of erected triangle is $-\tfrac{\sqrt3}{4}|Q-P|^2$ & \decl{Napoleon.apex_outward} & \src{napoleon}{Napoleon}{Basic}{112} & none (clockwise only if $P\ne Q$) \\
\multicolumn{4}{@{}l}{\emph{Lami (Section~\ref{sec:lami}); $V$ real inner product space, $\dim V=2$ and oriented where stated}}\\
Equilibrium $\Rightarrow$ equal cyclic area forms & \decl{Lami.areaForm_cyclic} & \src{lami}{Lami}{Basic}{27} & $\dim 2$, oriented; $F_1+F_2+F_3=0$ \\
$\omega(x,y)=\|x\|\|y\|\sin\measuredangle(x,y)$ & \decl{Lami.areaForm_eq_norm_mul_norm_mul_sin_oangle} & \src{lami}{Lami}{Basic}{36} & $\dim 2$, oriented; none (all $x,y$) \\
Lami, oriented product form & \decl{Lami.lami_oriented_mul} & \src{lami}{Lami}{Basic}{49} & $\dim 2$, oriented; $F_1+F_2+F_3=0$ \\
Lami, oriented ratio form & \decl{Lami.lami_oriented} & \src{lami}{Lami}{Basic}{100} & $\dim 2$, oriented; $F_1+F_2+F_3=0$; $F_1,F_2,F_3\ne0$ \\
Lami for forces in $\CC$ & \decl{Lami.lami_complex} & \src{lami}{Lami}{Basic}{111} & $V=\CC$; $F_1+F_2+F_3=0$; $F_1,F_2,F_3\ne0$ \\
Lami, unoriented product form & \decl{Lami.lami_mul} & \src{lami}{Lami}{Basic}{134} & any dimension; $F_1+F_2+F_3=0$ \\
Lami, unoriented ratio form & \decl{Lami.lami} & \src{lami}{Lami}{Basic}{148} & any dimension; $F_1+F_2+F_3=0$; $F_1,F_2,F_3\ne0$ \\
Lami with nonzero sines (Thm.~\ref{thm:lami}) & \decl{Lami.lami_nondegenerate} & \src{lami}{Lami}{Basic}{180} & any dimension; $F_1+F_2+F_3=0$; $F_1,F_2,F_3\ne0$; $\angle(F_1,F_2)\ne0,\pi$ \\
Parallel case: all sines $0$ & \decl{Lami.lami_degenerate} & \src{lami}{Lami}{Basic}{189} & any dimension; $F_1+F_2+F_3=0$; $F_1,F_2,F_3\ne0$; $\angle(F_1,F_2)\in\{0,\pi\}$ \\
Division-free form & \decl{Lami.lami_cross} & \src{lami}{Lami}{Basic}{211} & any dimension; $F_1+F_2+F_3=0$; $F_3\ne0$ \\
Converse, area-form version & \decl{Lami.equilibrium_of_areaForm} & \src{lami}{Lami}{Basic}{247} & $\dim 2$, oriented; $\omega(F_1,F_2)=\omega(F_2,F_3)=\omega(F_3,F_1)$; $\omega(F_1,F_2)\ne0$ \\
Converse, oriented ratio form & \decl{Lami.lami_oriented_converse} & \src{lami}{Lami}{Basic}{262} & $\dim 2$, oriented; $F_1,F_2,F_3\ne0$; three oriented sines $\ne0$; the two ratio equalities \\
\end{longtable}
}

\section{The method}\label{sec:method}

All five proofs follow the same three steps.
\begin{enumerate}
\item \emph{Coordinates.} Points are vectors of $K^3$ (homogeneous coordinates, Desargues),
pairs in $K\times K$ (Monge), complex numbers (Simson, Napoleon), or vectors of a real inner
product space (Lami). Lines, feet, apexes and centres of similitude are given by explicit
formulas in these coordinates.
\item \emph{Certificate.} The conclusion, after clearing denominators that the hypotheses make
nonzero, is an identity in the coordinates. Either it holds identically (Monge), or it is an
explicit multiple of a hypothesis polynomial (Napoleon, the concyclicity step of Simson), or it
factors so that one factor is forced to vanish by a hypothesis (Desargues, Simson). The identity
is closed by \lean{ring} (after \lean{field\_simp} where there are denominators) or by
\lean{linear\_combination} with the multiplier written in the source.
\item \emph{Translation.} The coordinate conclusion is turned into a standard Mathlib notion:
\mdecl{Projectivization.IsCollinear} for Desargues, \mdecl{Collinear} for Simson, \mdecl{dist}
for Napoleon, and norms with \mdecl{Orientation.oangle} or \mdecl{InnerProductGeometry.angle} for
Lami. Monge's theorem itself stops at the coordinate determinant; the bridge
\decl{Monge.det3_eq_zero_iff_collinear} identifies it with \mdecl{Collinear}, but the composed
statement is not stated as a declaration in the main files (Section~\ref{sec:limits}); the
Layer-14 file of version~2 states it (\decl{Monge.l14_monge}, Section~\ref{sec:l1415}).
\end{enumerate}
The certificates are small. The largest is Desargues's, of degree~12 in eighteen variables; it
takes about nine seconds to check. None of the proofs uses \lean{decide}, \lean{native\_decide} or
\lean{polyrith}; every multiplier given to \lean{linear\_combination} is written in the source.

The interest of the method is not that it is new (it is not; see~\cite{Chou}) but that it makes
the role of each hypothesis visible. A hypothesis is used in exactly one of three ways: to make a
denominator nonzero so that \lean{field\_simp} can clear it; to supply the polynomial on the right
of the certificate; or to make a translation lemma apply. Table~\ref{tab:nondeg} at the end of
Section~\ref{sec:lami} records which is which.

\section{Desargues's theorem}\label{sec:desargues}

\paragraph{Representation.} Let $K$ be a field. A point of $\PP^2(K)$ is represented by a nonzero
$v\in K^3$, and Mathlib's \mdecl{Projectivization.mk}~\lean{K v hv} is the point it represents. The
bracket \decl{Desargues.br}~$u\,v\,w=u\cdot(v\times w)$ is the determinant with rows $u,v,w$
(\decl{Desargues.br_eq_det}, from Mathlib's \mdecl{triple_product_eq_det}). If nonzero $u,v$
represent \emph{distinct points} of $\PP^2(K)$, that is, $u$ and $v$ are linearly independent,
then $u\times v\ne0$ represents the line through them. Distinctness of the representing vectors
is not enough: over $\QQ$, $u=(1,0,0)$ and $v=(2,0,0)$ satisfy $u\ne v$ but represent the same point, and
$u\times v=0$. (In characteristic~$2$ the vector $v$ is zero and represents no point at all.) Dually, if $l,m$ represent distinct lines then $l\times m\ne0$ represents their
intersection. For nonzero $o,a,a'$ the condition $[o\,a\,a']=0$ says that the three points
represented by $o,a,a'$ are collinear. When $a$ and $a'$ represent distinct points this is the
statement ``$O$ lies on the line $AA'$''; when they represent the same point, the line $AA'$ is
not defined and $[o\,a\,a']=0$ holds for every $o$. All perspective hypotheses below are stated
as $[o\,a\,a']=0$, which is meaningful in both cases.

\paragraph{The certificate.} For all $a,b,c,a',b',c'$ in $R^3$, $R$ any commutative ring,
\begin{equation}\label{eq:des}
\bigl[(a\times b)\times(a'\times b'),\ (b\times c)\times(b'\times c'),\ (c\times a)\times(c'\times a')\bigr]
 = [abc]\,[a'b'c']\,\bigl[a\times a',\ b\times b',\ c\times c'\bigr].
\end{equation}
This is \decl{Desargues.desargues_det_identity}. The proof rewrites each double cross product with
Mathlib's \mdecl{cross_cross_eq_smul_sub_smul} ($(u\times v)\times w=(u\cdot w)v-(v\cdot w)u$),
expands into coordinates, and closes the result with \lean{ring}. No hypothesis appears: the
identity is unconditional. The three points on the left are the Desargues points
$P=AB\cap A'B'$, $Q=BC\cap B'C'$, $R=CA\cap C'A'$, and the last factor on the right is the bracket
of the three joining lines $AA'$, $BB'$, $CC'$.

\paragraph{From the certificate to the theorem.} If a nonzero $o$ satisfies
$[o\,a\,a']=[o\,b\,b']=[o\,c\,c']=0$, then $o$ is orthogonal to the three vectors $a\times a'$, $b\times b'$, $c\times c'$, which are
therefore linearly dependent, so their bracket vanishes
(\decl{Desargues.br_eq_zero_of_dotProduct_eq_zero}, via Mathlib's
\mdecl{Matrix.exists_mulVec_eq_zero_iff}). By~\eqref{eq:des}, $[PQR]=0$.

\begin{theorem}[Desargues, determinant form]\label{thm:desdet}
Let $K$ be a field and $o,a,b,c,a',b',c'\in K^3$ with $o\ne0$. If $[o\,a\,a']=[o\,b\,b']=[o\,c\,c']=0$,
then $[PQR]=0$, where $P,Q,R$ are the cross-product expressions on the left of~\eqref{eq:des}.
\end{theorem}

\begin{lstlisting}
theorem desargues_det {o a b c a' b' c' : Fin 3 → K} (ho : o ≠ 0)
    (hOa : br o a a' = 0) (hOb : br o b b' = 0) (hOc : br o c c' = 0) :
    br ((a ⨯₃ b) ⨯₃ (a' ⨯₃ b')) ((b ⨯₃ c) ⨯₃ (b' ⨯₃ c')) ((c ⨯₃ a) ⨯₃ (c' ⨯₃ a')) = 0
\end{lstlisting}

The only non-degeneracy hypothesis is $o\ne0$, and it is needed: with $o=0$ the perspective
hypotheses are vacuous and the statement would claim that $[PQR]=0$ for all pairs of
triangles, which is false (the project's negative control found 2{,}946 of 2{,}985 random
non-perspective configurations over $\QQ$ with $[PQR]\ne0$). Everything else is allowed:
coincident vertices, degenerate triangles, coincident corresponding sides. In those cases one of
$P,Q,R$ may be the zero vector and the conclusion says nothing geometric, but it is still true.

\begin{theorem}[Desargues in $\PP^2(K)$]\label{thm:desmk}
Let $K$ be a field and let $A,B,C,A',B',C',O$ be points of $\PP^2(K)$. Suppose $A\ne B$, $B\ne C$,
$C\ne A$, $A'\ne B'$, $B'\ne C'$, $C'\ne A'$, and $AB\ne A'B'$, $BC\ne B'C'$, $CA\ne C'A'$ as
lines. Let $o,a,\dots,c'$ be nonzero representatives of $O,A,\dots,C'$, and suppose
$[o\,a\,a']=[o\,b\,b']=[o\,c\,c']=0$ (that is, $O,A,A'$ are collinear, and likewise for the other
two pairs; when $A\ne A'$ this says that $O$ lies on $AA'$). Then the points $AB\cap A'B'$,
$BC\cap B'C'$ and $CA\cap C'A'$ are collinear.
\end{theorem}

This is \decl{Desargues.desargues_mk}. Points are written \lean{mk K a ha} and so on; the line
through two points and the intersection of two lines are both Mathlib's
\mdecl{Projectivization.cross}; and the conclusion is \mdecl{Projectivization.IsCollinear} of the
three-point set. The perspective hypotheses are stated on representatives
(\lean{br o a a' = 0}); no hypothesis $A\ne A'$ is imposed, and a coincidence $A=A'$ simply makes
the corresponding perspective hypothesis vacuous. The six distinctness hypotheses on vertices and the three on sides have a
precise role. Mathlib's \mdecl{Projectivization.cross} returns its first argument when the cross
product of representatives vanishes, so it computes the line through $X$ and $Y$ only when
$X\ne Y$, and the intersection of two lines only when they are distinct
(\mdecl{Projectivization.cross_mk_of_ne}). The nine hypotheses are exactly what is needed to
rewrite every \lean{cross} in the statement as the cross product of representatives, after which
Theorem~\ref{thm:desdet} and \decl{Desargues.isCollinear_mk_of_br_eq_zero} finish the proof. They
are not needed for~\eqref{eq:des}. No hypothesis is placed on $O$ beyond its being a point; in
particular $O$ may coincide with a vertex, and the triangles need not be non-degenerate. No
assumption on the characteristic of~$K$ is needed.

\paragraph{The converse.} If both triangles are non-degenerate, $[abc]\ne0\ne[a'b'c']$, then by
\eqref{eq:des} $[PQR]=0$ forces $[AA'\,BB'\,CC']=0$
(\decl{Desargues.desargues_converse_det}). Non-degeneracy is necessary for this argument: if
either triangle is degenerate, the right-hand side of~\eqref{eq:des} vanishes for every
configuration and gives no information. The projective form
\decl{Desargues.desargues_converse_mk} assumes that the joining vectors $a\times a'$, $b\times b'$,
$c\times c'$ and the Desargues points $P,Q,R$ are nonzero, and concludes that the three lines
$AA'$, $BB'$, $CC'$, regarded as points of the dual plane, are collinear. That this is concurrency
of the three lines in the primal plane is the content of
\decl{Desargues.exists_dotProduct_eq_zero_of_br_eq_zero}: a vanishing bracket gives a nonzero $o$
orthogonal to all three, that is, a point on all three lines. A reader looking for a statement of
the form ``there is a point $O$ on $AA'$, $BB'$ and $CC'$'' has to compose the two declarations;
Mathlib has no concurrency predicate. Version~2 does this composition in the Layer-15 file:
\decl{Desargues.Classical.cov_A1_b} gives a point on the three lines, with lines as spans in
Mathlib's projective plane (Section~\ref{sec:l1415}).

\paragraph{Non-vacuity.} \path{Examples.lean} works out the configuration $O=(0{:}0{:}1)$,
$A=(1{:}0{:}0)$, $B=(0{:}1{:}0)$, $C=(1{:}1{:}1)$, $A'=(2{:}0{:}3)$, $B'=(0{:}5{:}{-1})$,
$C'=(4{:}4{:}11)$ over $\QQ$. Both triangles are non-degenerate (\decl{Desargues.ex_nonvacuous}),
the triangles are in perspective from~$O$ (\decl{Desargues.ex_perspective}), the Desargues points
are $P=(-2{:}-15{:}0)$, $Q=(-4{:}-39{:}-4)$, $R=(2{:}-12{:}-12)$, all nonzero
(\decl{Desargues.ex_points_ne_zero}), and $[PQR]=0$ is obtained both from the theorem
(\decl{Desargues.ex_desargues}) and by direct computation (\decl{Desargues.ex_desargues_compute}).

\section{Monge's theorem}\label{sec:monge}

For circles with centres $O_1,O_2\in K^2$ and radii $r_1\ne r_2$, the external centre of
similitude is the point dividing $O_1O_2$ externally in the ratio $r_1:r_2$,
\[
S_{12}=\frac{r_1O_2-r_2O_1}{r_1-r_2},
\]
which is \decl{Monge.ext}~\lean{O₁ O₂ r₁ r₂}, defined componentwise. Collinearity is the vanishing
of \decl{Monge.det3}~$A\,B\,C=(B-A)\times(C-A)$, the signed double area; over every field,
$\det_3(A,B,C)=0$ is equivalent to Mathlib's \lean{Collinear K \{A, B, C\}} in $K\times K$
(\decl{Monge.det3_eq_zero_iff_collinear}).

\begin{theorem}[Monge]\label{thm:monge}
Let $K$ be a field, $O_1,O_2,O_3\in K^2$, and $r_1,r_2,r_3\in K$ pairwise distinct. Then
$\det_3(S_{12},S_{23},S_{31})=0$.
\end{theorem}

\begin{lstlisting}
theorem monge (O₁ O₂ O₃ : K × K) (r₁ r₂ r₃ : K)
    (h12 : r₁ ≠ r₂) (h23 : r₂ ≠ r₃) (h31 : r₃ ≠ r₁) :
    det3 (ext O₁ O₂ r₁ r₂) (ext O₂ O₃ r₂ r₃) (ext O₃ O₁ r₃ r₁) = 0
\end{lstlisting}

\paragraph{The certificate.} After multiplying by $(r_1-r_2)(r_2-r_3)(r_3-r_1)$ the determinant
is a polynomial in the nine variables that vanishes identically. The proof is
\lean{unfold det3 ext; field\_simp; ring}. There is no hypothesis polynomial on the right: the
three hypotheses are used only to make the three denominators nonzero.

\paragraph{Hypotheses.} \emph{Scope of the geometric discussion.} Everything said below about
circles, centres of similitude, homotheties and tangents concerns genuine circles, with
$K=\RR$ and positive real radii. It is not valid for every field element that
\decl{Monge.monge} admits. For instance, $r_1=1$, $r_2=2$, $r_3=0$ satisfies the three
hypotheses, but then (as for any $r_1,r_2\ne0$ with $r_3=0$) $S_{23}=S_{31}=O_3$, so two of the three points coincide, the line
$S_{23}S_{31}$ is undefined, and the conclusion $\det_3=0$ holds trivially. More generally, if
$r_i=0$ the homothety ratio $r_j/r_i$ displayed below is undefined classically (Lean would read it
as~$0$). The Lean theorem is an identity over an arbitrary field; its geometric reading is
restricted to $r_1,r_2,r_3>0$ pairwise distinct.

Distinct radii are exactly the condition that each $S_{ij}$ is a finite
point. They cannot simply be dropped from the Lean statement: Lean defines $x/0=0$, so with
$r_1=r_2$ the term \lean{ext O₁ O₂ r₁ r₂} is the origin, and the determinant need not vanish. For
example, with centres $(0,0),(4,0),(0,3)$ and radii $1,1,3$, a computation in exact rational
arithmetic with Lean's division convention gives $\det_3=-9$. (This computation was done in Python,
not in Lean.) Geometrically, the equal-radius cases are of several kinds, none of which is
formalised. If exactly two radii agree, say $r_1=r_2\ne r_3$, and $O_1\ne O_2$, then $S_{12}$ is
the point at infinity in the direction of $O_1O_2$, and the projective form of Monge's theorem
says that $S_{23}S_{31}$ is parallel to $O_1O_2$ (a homothety centred at $O_3$ maps $O_1,O_2$ to
$S_{31},S_{23}$). If all three radii agree and the centres are distinct, then \emph{all three}
centres of similitude are at infinity and the statement degenerates to the triviality that three
points of the line at infinity are collinear; there is no finite Monge line. If two circles are
identical ($O_i=O_j$, $r_i=r_j$), every point of the plane is fixed by the identity map and there
is no distinguished centre of similitude at all. No hypothesis on the centres is needed for the
\emph{identity}: it holds for all $O_i$, including coincident centres. For $r_i\ne r_j$ and
$O_i=O_j$ (concentric circles) the formula gives $S_{ij}=O_i$, which is indeed the centre of the
homothety $x\mapsto O_i+\tfrac{r_j}{r_i}(x-O_i)$, but the circles have no common tangents.
More generally, the formula is meaningful without any geometric configuration assumption only as
the centre of the positive-ratio homothety; its classical description as the meeting point of the
two external common tangents requires that those tangents exist, which fails when one circle
lies strictly inside the other.
Nor are the radii required to be positive in the formal statement. Because \decl{Monge.monge} is universally quantified
over $r_1,r_2,r_3\in K$, replacing $r_2$ by $-r_2$ turns $S_{12}$ and $S_{23}$ into the
\emph{internal} centres of similitude, and gives the classical variant ``two internal centres and
the third external centre are collinear'' under the hypotheses $r_1\ne-r_2$, $-r_2\ne r_3$,
$r_3\ne r_1$, where now $r_1,r_2,r_3>0$ are the actual radii. We note this only as an instantiation; it is not stated as a separate declaration in the main files.
Version~2 proves the companion form (one external and two internal centres, for circles given as
Mathlib spheres) by the same sign-change argument, using $r_3\mapsto-r_3$ to obtain $E_{12},I_{23},I_{13}$, as \decl{Monge.Classical.cov_A2_b}
(Section~\ref{sec:l1415}).

\paragraph{Non-vacuity.} For centres $(0,0),(4,0),(0,3)$ and radii $1,2,3$ over $\QQ$, the
external centres are $(-4,0)$, $(12,-6)$ and $(0,-3/2)$ (\decl{Monge.monge_example_values}), and
$\det_3=0$ (\decl{Monge.monge_example}).

\paragraph{Scope.} Neither circles nor tangent lines appear in the formal statement. The radii are
elements of $K$ and \decl{Monge.ext} is a formula. The bridge \decl{Monge.ext_eq_lineMap_homothety}
shows, for $r_1\ne0$ and $r_1\ne r_2$, that $S_{12}$ is the point
\lean{AffineMap.lineMap O₁ O₂ (r₁ / (r₁ - r₂))} of the line $O_1O_2$ and that the homothety
(\mdecl{AffineMap.homothety}) with centre $S_{12}$ and ratio $r_2/r_1$ sends $O_1$ to $O_2$. The
hypothesis $r_1\ne0$ is needed in addition to $r_1\ne r_2$: for $r_1=0$ the formula gives
$S_{12}=O_1$, which is not a centre of similitude, and the ratio $r_2/r_1$ is undefined (Lean reads
it as~$0$). The lemma is about centres only; in the main files, that this homothety maps the one circle onto the other
is not formalised. Version~2 proves it in the Layer-15 file: for two circles with distinct centres and positive
unequal radii given as Mathlib spheres, a point at which the homothety of ratio $r_t/r_s$ maps the
one circle onto the other exists and is unique (\decl{Monge.Classical.cov_A2_obligation_ext}), it is
the point \decl{Monge.ext} (\decl{Monge.Classical.ext_eq_of_isExt}), and for three circles with
pairwise distinct radii the three such points are \mdecl{Collinear}
(\decl{Monge.Classical.cov_A2_a}). Not formalised is the formulation in which the external common
tangents of each pair meet at~$S_{ij}$ (a formulation that applies only when those tangents
exist, so not to nested circles). Mathlib's \path{Geometry/Euclidean/MongePoint.lean}
concerns the Monge point of a simplex, a different theorem, and is not used.

\paragraph{Prior formalisations.} This theorem is not new to Lean, and our proof is not the first;
two earlier complete Lean proofs exist (Section~\ref{sec:related}).

\section{The Simson line}\label{sec:simson}

\paragraph{Representation.} The Euclidean plane is $\CC$, regarded as a real affine space. For
$u\ne v$ the foot of the perpendicular from $p$ to the line $uv$ is
\decl{Simson.perpFoot}~$p\,u\,v=u+t(v-u)$ with
$t=\langle p-u,v-u\rangle/|v-u|^2$, written in real and imaginary parts
(\decl{Simson.footParam}). Two lemmas say that this is the foot: it lies on Mathlib's
\lean{line[ℝ, u, v]} (\decl{Simson.perpFoot_mem_line}), and $p-\mathrm{foot}$ is orthogonal to
$v-u$ when $u\ne v$ (\decl{Simson.inner_perpFoot}, stated in coordinates). Moreover the parameter
is $t=\langle p-u,v-u\rangle/\|v-u\|^2$ for Mathlib's real inner product on $\CC$
(\decl{Simson.footParam_eq_inner}), and the foot is Mathlib's
\mdecl{EuclideanGeometry.orthogonalProjection} of $p$ onto \lean{line[ℝ, u, v]}
(\decl{Simson.perpFoot_eq_orthogonalProjection}), with no hypothesis (for $u=v$ both sides
are~$u$). The signed double area is
\decl{Simson.area2}, and \decl{Simson.concyc}~$a\,b\,c\,p$ is the $4\times4$ determinant with rows
$(x,y,x^2+y^2,1)$, written out by cofactor expansion. For a non-degenerate triangle,
$\mathrm{area2}(a,b,c)\ne0$, it vanishes exactly when $\{a,b,c,p\}$ is Mathlib's
\mdecl{EuclideanGeometry.Cospherical} (\decl{Simson.concyc_eq_zero_iff_cospherical}). The
hypothesis cannot be dropped, and $\mathrm{concyc}$ does not vanish merely because $a,b,c$ are
collinear: for $a,b,c=(0,0),(1,0),(2,0)$ and $p=(0,1)$ one has $\mathrm{concyc}(a,b,c,p)=2$, while
four points on one line have $\mathrm{concyc}=0$ without being cospherical. (The comment on
\decl{Simson.concyc} in the cited source describes its zero set inaccurately; the bridge lemma,
with its hypothesis, is the correct statement.)

\paragraph{The certificate.} Write $D_1=|c-b|^2$, $D_2=|a-c|^2$, $D_3=|b-a|^2$ and $S=\mathrm{area2}(a,b,c)$.
For pairwise distinct $a,b,c$ and every $p$,
\begin{equation}\label{eq:simson}
\mathrm{area2}\bigl(\mathrm{foot}_{bc}(p),\mathrm{foot}_{ca}(p),\mathrm{foot}_{ab}(p)\bigr)\cdot D_1D_2D_3
 = S^2\cdot\mathrm{concyc}(a,b,c,p).
\end{equation}
This is \decl{Simson.area2_perpFoot_mul}, proved by \lean{field\_simp; ring}; the distinctness
hypotheses make the $D_i$ nonzero (\decl{Simson.sqD_ne_zero}). A second, linear certificate shows
that four points at distance $r$ from a common $o$ have $\mathrm{concyc}=0$
(\decl{Simson.concyc_eq_zero_of_dist}): $\mathrm{concyc}$ is a combination of the four equations
$|x-o|^2=r^2$ with the cofactors as multipliers, given explicitly to \lean{linear\_combination}.
The factorisation \eqref{eq:simson} was found with \texttt{sympy} (\path{code/cert.py},
\path{code/gen.py}, \path{code/gen2.py}) and then checked by Lean.

\begin{theorem}[Wallace--Simson, forward direction]\label{thm:simson}
Let $a,b,c\in\CC$ be pairwise distinct and let $p\in\CC$ be such that $\{a,b,c,p\}$ is
cospherical. Then the feet of the perpendiculars from $p$ to the lines $bc$, $ca$, $ab$ are
collinear.
\end{theorem}

\begin{lstlisting}
theorem simson {a b c p : ℂ} (hab : a ≠ b) (hbc : b ≠ c) (hca : c ≠ a)
    (hcirc : EuclideanGeometry.Cospherical ({a, b, c, p} : Set ℂ)) :
    Collinear ℝ ({perpFoot p b c, perpFoot p c a, perpFoot p a b} : Set ℂ)
\end{lstlisting}

The hypothesis is Mathlib's \mdecl{EuclideanGeometry.Cospherical}, which in the plane $\CC$ means
that the four points lie on a common circle (of any radius, including radius~0, which the
distinctness hypotheses exclude). The conclusion is Mathlib's \mdecl{Collinear}~$\RR$. The
intermediate \decl{Simson.simson_area} states the conclusion as $\mathrm{area2}=0$ with the
circle given by a centre $o$ and radius $r$.

\paragraph{Hypotheses.} The only hypotheses are distinctness of the three vertices and the circle
condition. Distinctness is used only to make the feet well defined and the $D_i$ nonzero. No
separate hypothesis that the triangle is non-degenerate is needed: three distinct points on a
circle are not collinear, and in any case \eqref{eq:simson} does not need it. The point $p$ may be
a vertex, and right-angled triangles are included; the identity covers all of these.

\paragraph{The converse.} The main files of our development prove only the forward direction
of the Wallace--Simson theorem. For a non-degenerate triangle ($S\ne0$, which implies
distinctness), \eqref{eq:simson} gives
$\mathrm{area2}(\text{feet})=0\iff\mathrm{concyc}(a,b,c,p)=0$
(\decl{Simson.area2_perpFoot_eq_zero_iff}). This is an equivalence between two \emph{local
determinant expressions}; it is not the geometric converse. Neither side is a Mathlib notion. On
the circle side the bridge is now two-way for $S\ne0$ (\decl{Simson.concyc_eq_zero_iff_cospherical}),
but on the feet side it is one-way: \decl{Simson.collinear_of_area2_eq_zero} turns
$\mathrm{area2}=0$ into \mdecl{Collinear}, and there is no lemma from \mdecl{Collinear} back to
$\mathrm{area2}=0$. The geometric converse, ``if the feet are collinear then $p$ lies on the
circumcircle'', was therefore not formalised in the first version of this paper; the missing step
was that converse bridge for $\mathrm{area2}$. Version~2 supplies it in the Layer-15 file:
\decl{Simson.Classical.area2_eq_zero_of_collinear} (eight lines) proves that three collinear points
have $\mathrm{area2}=0$, and \decl{Simson.Classical.cov_A3_b} chains it with the two lemmas above to
prove the converse for a Mathlib \mdecl{Affine.Triangle} in \lean{EuclideanSpace ℝ (Fin 2)}, with
conclusion membership in Mathlib's \mdecl{Affine.Simplex.circumsphere} (Section~\ref{sec:l1415}).
The converse was proved earlier in the Isabelle AFP entry and in
Guilhot's Rocq development cited in Section~\ref{sec:related}.

\paragraph{Non-vacuity.} The library contains an \lean{example} applying \decl{Simson.simson} to the
triangle $1,i,-1$ on the unit circle with $p=-i$, and a negative control showing that for $p=0$
(the centre) the pedal triangle has nonzero area.

\paragraph{Where the reals enter.} The certificate~\eqref{eq:simson} is an identity of
polynomials in the eight real coordinates $(x,y)$ of $a,b,c,p$ and, as a polynomial identity, it
holds over any commutative ring. The proof uses the reals in two places: the foot is defined
through the Euclidean inner product, and $D_i\ne0$ for $a\ne b$ uses that a sum of two real
squares vanishes only if both do. The second step is what fails if the coordinates are allowed to
range over a field such as $\CC$ (or any field containing a square root of $-1$) while $D_i$ is
kept as the \emph{bilinear} expression $(\Delta x)^2+(\Delta y)^2$, which is what the polynomial
identity contains: this quadratic form is isotropic there, e.g.\ $\Delta=(1,i)$ gives
$1^2+i^2=0$ although $\Delta\ne0$, so two distinct points can have $D=0$ and the foot is then
undefined. (The Hermitian expression $|\Delta x|^2+|\Delta y|^2$ is positive for $\Delta\ne0$, but
it is not a polynomial in the coordinates and does not appear in~\eqref{eq:simson}.) Over an
arbitrary field one would have to assume the side vectors non-isotropic. This is why the theorem
is stated in $\CC$ regarded as the real Euclidean plane, with real coordinates, and not over an
arbitrary field.

\section{Napoleon's theorem}\label{sec:napoleon}

\paragraph{Representation.} The plane is $\CC$. For $\omega\in\CC$, \decl{Napoleon.apex}~$\omega\,P\,Q
=P+\omega(Q-P)$ is the image of $Q$ under the map $z\mapsto P+\omega(z-P)$. For general $\omega$
this is a spiral similarity (or, for $\omega=0$, constant), not a rotation; it is a rotation about
$P$ exactly when $|\omega|=1$.
The two rotations are $\omega_o=\tfrac12-\tfrac{\sqrt3}2 i$ (by $-\pi/3$, \decl{Napoleon.ωo}) and
$\omega_i=\overline{\omega_o}$ (\decl{Napoleon.ωi}). The Napoleon points $G_a,G_b,G_c$ are the
centroids of the triangles erected on $BC$, $CA$, $AB$ (\decl{Napoleon.Ga} and its two
companions).

\paragraph{The certificate.} If $\omega^2-\omega+1=0$, then for all $A,B,C$
\begin{equation}\label{eq:nap}
G_b-G_a-\omega(G_c-G_a)=\tfrac13(A-2B+C)\,(\omega^2-\omega+1).
\end{equation}
This is \decl{Napoleon.key}, proved by \lean{linear\_combination (A - 2 * B + C) / 3 * h}. It says that
$G_b$ is the image of $G_c$ under $z\mapsto G_a+\omega(z-G_a)$. Over $\CC$ the hypothesis
$\omega^2-\omega+1=0$ already forces $\omega=e^{\pm i\pi/3}$, so this map is in fact a rotation by
$\pm\pi/3$; but \decl{Napoleon.key} is a ring identity and does not use that.

More precisely, \eqref{eq:nap} needs no hypothesis at all: both sides are polynomials in
$A,B,C,\omega$, and the identity holds for every $\omega$, not merely for roots of
$\omega^2-\omega+1$. The Lean declaration \decl{Napoleon.key} states the specialisation
$G_b-G_a=\omega(G_c-G_a)$ that is used below, and carries $\omega^2-\omega+1=0$ as a hypothesis;
as formalised it is therefore marginally weaker than the displayed identity, which we have kept in
the conditional form in which the geometry consumes it.

\begin{theorem}[Napoleon]\label{thm:napoleon}
For all $A,B,C\in\CC$, the centres $G_a,G_b,G_c$ of the equilateral triangles erected on the
sides by the rotation $\omega_o$ satisfy $|G_aG_b|=|G_bG_c|=|G_cG_a|$; and the same holds for
the rotation~$\omega_i$.
\end{theorem}

\begin{lstlisting}
theorem napoleon_outer (A B C : ℂ) :
    dist (Ga ωo A B C) (Gb ωo A B C) = dist (Gb ωo A B C) (Gc ωo A B C) ∧
    dist (Gb ωo A B C) (Gc ωo A B C) = dist (Gc ωo A B C) (Ga ωo A B C)
\end{lstlisting}

Both statements (\decl{Napoleon.napoleon_outer}, \decl{Napoleon.napoleon_inner}) are instances of
\decl{Napoleon.napoleon_general}, which holds for every $\omega$ with $\omega^2-\omega+1=0$,
$|\omega|=1$ and $|1-\omega|=1$. These three hypotheses are not independent: the roots of
$\omega^2-\omega+1$ are $\omega=\tfrac12\pm\tfrac{\sqrt3}2i=e^{\pm i\pi/3}$, and both satisfy
$|\omega|=|1-\omega|=1$, so the two norm hypotheses follow from the first. (Conversely,
$|\omega|=|1-\omega|=1$ alone also characterises $e^{\pm i\pi/3}$.) The Lean statement keeps all
three because the proof uses them in different places: the first drives \eqref{eq:nap}, and the
other two turn it into equalities of distances, since $G_b-G_a=\omega(G_c-G_a)$ and
$G_c-G_b=(1-\omega)(G_c-G_a)$. Deriving the norm hypotheses from the polynomial one would remove
two redundant binders; we have not done so. The lemma
\decl{Napoleon.apex_equilateral} checks that the erected triangles are equilateral, under the
same two norm hypotheses.

\paragraph{Hypotheses.} Theorem~\ref{thm:napoleon} has none. For a degenerate triangle the
conclusion still holds, possibly with a Napoleon triangle that is a single point. Orientation
enters only in the interpretation of $\omega_o$ as ``outward''. The lemma
\decl{Napoleon.apex_outward} computes the signed area
$\mathrm{signedArea}(P,Q,\mathrm{apex}\,\omega_o\,P\,Q)=-\tfrac{\sqrt3}4|Q-P|^2$, for all $P,Q$. So
every triangle erected with $\omega_o$ on a base with $P\ne Q$ is clockwise; when $P=Q$ the erected
``triangle'' is the single point $P$, its signed area is~$0$, and it has no orientation. If $ABC$
is counter-clockwise ($\mathrm{signedArea}(A,B,C)>0$, which forces the vertices to be distinct),
each apex therefore lies on the other side of its base from the third vertex, and the $\omega_o$
triangles are the outer ones; for a clockwise $ABC$ the outer triangles are the $\omega_i$ ones,
and \decl{Napoleon.napoleon_inner} covers that case. ``Outward'' is thus expressed through signed
area rather than through half-planes; in the main files the step from the signed-area statement to ``on the opposite
side of the line'' is not formalised. Version~2 adds it: \decl{Napoleon.l14_napoleon_outer} takes
apexes defined by equal sides and Mathlib's strict opposite side
(\mdecl{AffineSubspace.SOppSide}), with no orientation hypothesis, and
\decl{Napoleon.Classical.cov_A4_a} proves the classical outer statement in
\lean{EuclideanSpace ℝ (Fin 2)}, including that the outer Napoleon triangle of a triangle is
non-degenerate (Section~\ref{sec:l1415}). Mixing orientations (erecting some triangles with $\omega_o$
and some with $\omega_i$) makes the statement false; the project's numerical check found a spread
of 1.15 between the three distances in such a configuration.

\paragraph{Non-vacuity.} An \lean{example} checks that for $A=0$, $B=1$, $C=i$ the triangle is
counter-clockwise, the Napoleon equality holds, and $|G_aG_b|\ne0$.

\paragraph{Prior formalisations.} This theorem is not new to Lean, and our version is not the first
in a proof assistant to treat the exterior condition; see Section~\ref{sec:related}.

\section{Lami's theorem}\label{sec:lami}

Lami's theorem says that if three forces $F_1,F_2,F_3$ acting at a point are in equilibrium,
$F_1+F_2+F_3=0$, then
\[
\frac{|F_1|}{\sin\alpha}=\frac{|F_2|}{\sin\beta}=\frac{|F_3|}{\sin\gamma},
\]
where $\alpha$ is the angle between $F_2$ and $F_3$, $\beta$ the angle between $F_3$ and $F_1$, and
$\gamma$ the angle between $F_1$ and~$F_2$.

\paragraph{Oriented form and the certificate.} Let $V$ be a real inner product space of dimension
two with an orientation $o$, and let $\omega$ be its area form (Mathlib's
\mdecl{Orientation.areaForm}). Substituting $F_3=-F_1-F_2$ and using bilinearity and
$\omega(x,x)=0$ gives
\begin{equation}\label{eq:lami}
\omega(F_1,F_2)=\omega(F_2,F_3)=\omega(F_3,F_1)
\end{equation}
(\decl{Lami.areaForm_cyclic}, closed by \lean{ring} after \lean{simp}). This is the certificate.
The translation to sines is \decl{Lami.areaForm_eq_norm_mul_norm_mul_sin_oangle},
$\omega(x,y)=\|x\|\,\|y\|\sin\measuredangle(x,y)$ for all $x,y$ including zero, proved through
Mathlib's K\"ahler form and \mdecl{Complex.sin_arg}. Together they give the product form
\decl{Lami.lami_oriented_mul}, with equilibrium as the only hypothesis:
$\|F_2\|\|F_3\|\sin\measuredangle(F_2,F_3)=\|F_3\|\|F_1\|\sin\measuredangle(F_3,F_1)=\|F_1\|\|F_2\|\sin\measuredangle(F_1,F_2)$.
The ratio form \decl{Lami.lami_oriented} needs $F_1,F_2,F_3\ne0$ and follows from the product
form by \decl{Lami.ratio_of_mul}: if all $n_i\ne0$ and
$n_2n_3s_1=n_3n_1s_2=n_1n_2s_3$, then either all $s_i$ vanish or none does, and
$n_1/s_1=n_2/s_2=n_3/s_3$ in Lean's convention. \decl{Lami.lami_complex} is the instance
$V=\CC$ with its standard orientation.

\paragraph{Unoriented form.} With Mathlib's unoriented angle \mdecl{InnerProductGeometry.angle}
(values in $[0,\pi]$), $\|x\|\|y\|\sin\angle(x,y)$ is the square root of the Gram determinant
$\langle x,x\rangle\langle y,y\rangle-\langle x,y\rangle^2$
(\decl{Lami.norm_mul_norm_mul_sin_angle}). The three Gram determinants of $(F_2,F_3)$,
$(F_3,F_1)$, $(F_1,F_2)$ agree once $F_3=-F_1-F_2$ is substituted, again by \lean{ring}. This
gives \decl{Lami.lami_mul} and \decl{Lami.lami} in every real inner product space; no
two-dimensionality is assumed, since the three forces span a space of dimension at most two.

\begin{theorem}[Lami, non-degenerate form]\label{thm:lami}
Let $V$ be a real inner product space and $F_1,F_2,F_3\in V$ nonzero with $F_1+F_2+F_3=0$ and
$\angle(F_1,F_2)\notin\{0,\pi\}$. Then $\sin\alpha$, $\sin\beta$, $\sin\gamma$ are nonzero and
$\|F_1\|/\sin\alpha=\|F_2\|/\sin\beta=\|F_3\|/\sin\gamma$.
\end{theorem}

\begin{lstlisting}
theorem lami_nondegenerate {F₁ F₂ F₃ : V} (h : F₁ + F₂ + F₃ = 0)
    (h₁ : F₁ ≠ 0) (h₂ : F₂ ≠ 0) (h₃ : F₃ ≠ 0)
    (hγ : angle F₁ F₂ ≠ 0) (hγ' : angle F₁ F₂ ≠ π) :
    (Real.sin (angle F₂ F₃) ≠ 0 ∧ Real.sin (angle F₃ F₁) ≠ 0 ∧ Real.sin (angle F₁ F₂) ≠ 0) ∧
      ‖F₁‖ / Real.sin (angle F₂ F₃) = ‖F₂‖ / Real.sin (angle F₃ F₁) ∧
      ‖F₂‖ / Real.sin (angle F₃ F₁) = ‖F₃‖ / Real.sin (angle F₁ F₂)
\end{lstlisting}

This is the declaration to cite for the classical theorem. The statement \decl{Lami.lami} has the
same conclusion without the sine conjunct and without the non-parallel hypotheses; it is true, but
in the parallel case it is true only because Lean defines $x/0=0$. \decl{Lami.lami_degenerate}
makes this explicit: if $F_1\parallel F_2$ then all three sines vanish. The division-free form
\decl{Lami.lami_cross}, $\|F_1\|\sin\beta=\|F_2\|\sin\alpha$, needs only $F_3\ne0$.

\paragraph{Converse.} In dimension two, if the three cyclic area forms are equal and nonzero, the
forces are in equilibrium (\decl{Lami.equilibrium_of_areaForm}, via Mathlib's
\mdecl{Orientation.inner_mul_areaForm_sub}). Hence three nonzero forces with nonzero oriented sines
whose magnitudes are proportional to the sines of the opposite oriented angles are in equilibrium
(\decl{Lami.lami_oriented_converse}). Orientation is essential here: with unoriented angles the
converse is false, because replacing $F_1$ by $-F_1$ leaves every unoriented sine unchanged but
destroys equilibrium.

\paragraph{The law of sines.} Lami's theorem is the law of sines for the force triangle with sides
$\|F_1\|,\|F_2\|,\|F_3\|$ and interior angles $\pi-\alpha,\pi-\beta,\pi-\gamma$. The declaration
\decl{Lami.lami_via_law_sin} only restates Mathlib's
\mdecl{InnerProductGeometry.sin_angle_mul_norm_eq_sin_angle_mul_norm}. It is not used in the
proofs, and the connection between Lami's theorem and the force triangle is not formalised.

{\small
\begin{table}[t]
\centering
\begin{tabular}{@{}>{\raggedright\arraybackslash}p{0.13\textwidth}
  >{\raggedright\arraybackslash}p{0.24\textwidth}
  >{\raggedright\arraybackslash}p{0.27\textwidth}
  >{\raggedright\arraybackslash}p{0.28\textwidth}@{}}
\toprule
Theorem & Hypothesis & Role & Without it \\ \midrule
Desargues & $o\ne0$ & turns perspective into $[AA'BB'CC']=0$ & statement false \\
 & $[o\,a\,a']=[o\,b\,b']=[o\,c\,c']=0$ (forward) & perspective from $O$; with $o\ne0$ gives $[(a{\times}a')(b{\times}b')(c{\times}c')]=0$, the third factor of~\eqref{eq:des} & no theorem (the hypothesis) \\
 & six vertex inequalities, three side inequalities, nonzero representatives & translation: Mathlib's \lean{cross} computes lines and meets & conclusion about the wrong points \\
 & $[PQR]=0$, $P,Q,R\ne0$, joins $\ne0$ (converse) & hypothesis of the converse; translation & --- \\
 & $[abc],[a'b'c']\ne0$ (converse) & cancels a factor of~\eqref{eq:des} & no information \\
Monge & $r_i\ne r_j$ & clears denominators & false in Lean ($x/0=0$); point at infinity \\
 & (centres) & not needed & --- \\
 & $r_1\ne0$, $r_1\ne r_2$ (bridge \decl{Monge.ext_eq_lineMap_homothety}) & clears $r_1$ and $r_1-r_2$ & for $r_1=0$, $S_{12}=O_1$ is not a centre of similitude \\
Simson & $a,b,c$ distinct & clears $D_i$; feet well defined & feet undefined (Lean: $x/0=0$) \\
 & cospherical & supplies $\mathrm{concyc}=0$ in~\eqref{eq:simson} & false (negative control) \\
 & $S\ne0$ (determinant equivalence) & cancels $S^2$ & no information \\
 & $S\ne0$ (bridge \decl{Simson.concyc_eq_zero_iff_cospherical}) & circumcentre formula has denominator $2S$ & false: four distinct collinear points have $\mathrm{concyc}=0$ \\
Napoleon & $\omega^2-\omega+1=0$ & multiplier in~\eqref{eq:nap} & still implied by the two retained norm hypotheses (not independently removable); false for $\omega$ violating all three \\
 & $|\omega|=|1-\omega|=1$ & translation to distances & redundant given $\omega^2-\omega+1=0$ (each set implies the other: both say $\omega=e^{\pm i\pi/3}$) \\
 & (non-degenerate $ABC$) & not needed & --- \\
Lami & $F_1+F_2+F_3=0$ & substitution in~\eqref{eq:lami} & --- \\
 & $F_1,F_2,F_3\ne0$ & ratio forms (\decl{Lami.lami}, \decl{Lami.lami_oriented}, \decl{Lami.lami_nondegenerate}) & ratios meaningless \\
 & $F_3\ne0$ only & division-free form \decl{Lami.lami_cross}: cancels $\|F_3\|$ & product form \decl{Lami.lami_mul} still holds \\
 & $F_1\nparallel F_2$ & sines nonzero & ratios hold only via $x/0=0$ \\
 & (converse, area form) $\omega$'s equal, $\omega(F_1,F_2)\ne0$ & hypotheses of \decl{Lami.equilibrium_of_areaForm} & no information \\
 & (converse, ratio form) $F_i\ne0$, oriented sines $\ne0$, two ratio equalities & hypotheses of \decl{Lami.lami_oriented_converse}; sines clear denominators & false with unoriented angles \\
\bottomrule
\end{tabular}
\caption{Non-degeneracy and other hypotheses of the five theorems, and how each is used.}
\label{tab:nondeg}
\end{table}
}

\section{Does the statement mean what it says?}\label{sec:correct}

Each development introduces some local notions. A reader has to check that they mean what their
names suggest. The lemmas available for this are as follows.
\begin{itemize}
\item \emph{Desargues.} The bracket is the determinant (\decl{Desargues.br_eq_det}); a vanishing
bracket is equivalent to Mathlib's \mdecl{Projectivization.IsCollinear} of the three points
(\decl{Desargues.isCollinear_mk_of_br_eq_zero}, \decl{Desargues.br_eq_zero_of_isCollinear_mk});
and a vanishing bracket of three lines is concurrency
(\decl{Desargues.exists_dotProduct_eq_zero_of_br_eq_zero},
\decl{Desargues.br_eq_zero_of_dotProduct_eq_zero}). Lines and intersections in
Theorem~\ref{thm:desmk} are Mathlib's own. For the worked example, the representatives are nonzero
(\decl{Desargues.ex_vertices_ne_zero}), the centre $O$ is Mathlib's intersection
\mdecl{Projectivization.cross} of the lines $AA'$ and $BB'$ (\decl{Desargues.exO_eq_cross}), and
each of $A',B',C'$ lies on the line from $O$ to the corresponding vertex in the sense of
\mdecl{Projectivization.IsCollinear} (\decl{Desargues.exA'_isCollinear},
\decl{Desargues.exB'_isCollinear}, \decl{Desargues.exC'_isCollinear}).
\item \emph{Monge.} \decl{Monge.det3}~$A\,B\,C=0$ is equivalent to Mathlib's \mdecl{Collinear} over
every field (\decl{Monge.det3_eq_zero_iff_collinear}). For $r_1\ne0$ and $r_1\ne r_2$,
\decl{Monge.ext} lies on the line $O_1O_2$ (\mdecl{AffineMap.lineMap}) and is the centre of the
homothety (\mdecl{AffineMap.homothety}) of ratio $r_2/r_1$ sending $O_1$ to $O_2$
(\decl{Monge.ext_eq_lineMap_homothety}); for $r_1=0$ it is not a centre of similitude
(Section~\ref{sec:monge}).
\item \emph{Simson.} \decl{Simson.perpFoot} is Mathlib's
\mdecl{EuclideanGeometry.orthogonalProjection} onto the line
(\decl{Simson.perpFoot_eq_orthogonalProjection}), and \decl{Simson.footParam} is the usual
inner-product coefficient (\decl{Simson.footParam_eq_inner}). \decl{Simson.area2} has a one-way
bridge in the main files, \decl{Simson.collinear_of_area2_eq_zero} ($\mathrm{area2}=0$ implies \mdecl{Collinear}), but
no lemma in the other direction there; version~2 adds it in the Layer-15 file
(\decl{Simson.Classical.area2_eq_zero_of_collinear}). \decl{Simson.concyc} is equivalent to \mdecl{Cospherical} under
the hypothesis $\mathrm{area2}(a,b,c)\ne0$ (\decl{Simson.concyc_eq_zero_iff_cospherical}), which
cannot be dropped (Section~\ref{sec:simson}). The hypothesis and the conclusion of
Theorem~\ref{thm:simson} are Mathlib notions.
\item \emph{Napoleon.} The conclusion uses Mathlib's \mdecl{dist}.
\decl{Napoleon.apex_equilateral} shows that the erected triangles are equilateral, and
\decl{Napoleon.apex_outward} fixes their orientation (for bases with $P\ne Q$). Each local
definition has a bridge: \decl{Napoleon.centroid} is Mathlib's \mdecl{Finset.centroid}
(\decl{Napoleon.centroid_eq_finset_centroid}), and so are the Napoleon points
(\decl{Napoleon.Ga_eq_centroid}, \decl{Napoleon.Gb_eq_centroid}, \decl{Napoleon.Gc_eq_centroid});
for a unit complex number, \decl{Napoleon.apex} is Mathlib's \mdecl{rotation} about~$P$
(\decl{Napoleon.apex_eq_rotation}); \decl{Napoleon.signedArea} is half of
\mdecl{Complex.orientation}'s area form (\decl{Napoleon.signedArea_eq_areaForm}); and
$\omega_o=e^{-i\pi/3}$, $\omega_i=e^{i\pi/3}$ (\decl{Napoleon.ωo_eq_exp}, \decl{Napoleon.ωi_eq_exp}).
\item \emph{Lami.} Hypotheses and conclusions use only norms and Mathlib's oriented and unoriented
angles. \decl{Lami.areaForm_eq_norm_mul_norm_mul_sin_oangle} and
\decl{Lami.norm_mul_norm_mul_sin_angle} relate the algebraic objects of the proofs to those
angles.
\end{itemize}
The house bridge inventory (Section~\ref{sec:verify}) lists every local definition and whether it
is bridged.

\section{Related work}\label{sec:related}

Each project began with a prior-art search of Mathlib (sources and the lists
\path{docs/100.yaml} and \path{docs/1000.yaml}), GitHub code, repository and pull-request search,
the Isabelle Archive of Formal Proofs~\cite{AFP}, and the web. The searches were made between
30~September and 3~October 2026. They are recorded in the \path{SPIKE_REPORT.md} of each
repository; we summarise them here.

\paragraph{Desargues.} Desargues's theorem is number~87 on Wiedijk's list of 100
theorems~\cite{Wiedijk100}, which records formalisations in HOL~Light (Harrison), Isabelle (Bordg),
Rocq (Guilhot and Narboux), Metamath (Megill) and Mizar (Leo\'nczuk and Pra\.zmowski). The
HOL~Light proof~\cite{HOLDesargues} is a close precedent for ours: it builds the real projective
plane from nonzero vectors of $\RR^3$, proves that collinearity of three points is the vanishing
of their bracket (\lean{COLLINEAR\_BRACKET}), and proves Desargues's theorem
(\lean{DESARGUES\_DIRECT}, following Richter-Gebert's bracket proof) by reducing it to bracket
identities closed by the ring decision procedure \lean{REAL\_RING}, under seventeen
non-collinearity hypotheses. It is thus also a coordinate certificate proof, over~$\RR$ rather
than over an arbitrary field, and with a different factorisation from~\eqref{eq:des}. In Isabelle,
the AFP entry \emph{Projective Geometry}~\cite{AFPProjective} proves the theorem from an
axiomatisation of projective space via matroid rank; Magaud, Narboux and Schreck give a Coq proof
from incidence axioms~\cite{MNS}; these two are synthetic. In Mathlib, \path{docs/100.yaml} (at
the pinned revision) lists the theorem with no declaration, and
\path{Projectivization/Constructions.lean} supplies the cross-product API that we use but no
Desargues statement. The open, unmerged Mathlib pull request \#37456~\cite{PR37456} (created
31~March 2026; head commit \texttt{0900538c}) proves the affine \emph{parallel} form of
Desargues's theorem in an affine space. We therefore do not claim a first Lean proof of
Desargues's theorem: we found no Lean proof of the \emph{projective-plane} formulation, while an
unmerged pull request covers the affine (parallel) form.

\paragraph{Monge.} Monge's three-circle theorem is \emph{not} new to Lean, and we withdraw the
claim, made in an earlier draft of this paper, that ours is the first Lean proof. Two
independent, complete Lean proofs exist. (i)~The benchmark repository
\path{epfl-lara/icml-26-lean-challenges}~\cite{ShadowBench}, file
\path{autformalization/problems/geometry/L2/geo_gen_L2_005/ShadowBench/Source/Main.lean} at commit
\texttt{7541fcecef97cc124f38355a2cd2ba0ec8485dd5} (12~June 2026; 135 lines, no \lean{sorry}),
proves \lean{EuclideanGeometry.monges\_circle\_theorem} (line~73), with the external homothety
centres given explicitly. This predates our development (3~October 2026) by about four months,
and on it alone the claim of priority is withdrawn.
(ii)~The repository \path{facebookresearch/atlas-lean}~\cite{AtlasLean}, file
\path{MathlibExt/Geometry/Euclidean/MongeCircles.lean}, added by pull request~\#62 at commit
\texttt{39febc62dfe68f9651c72eedf7916a7a070141a8} (5~October 2026, 21:40~UTC; 83 lines, no
\lean{sorry}), proves \lean{monge\_three\_circles\_general}, stated in an arbitrary
real module with distinct weights as the only hypothesis, and \lean{monge\_three\_circles}. Our
development was made independently of it.
\emph{Dates and what they show.} Our Monge development dates from 3~October 2026 and was made
public at \path{github.com/jpbald93/monge} on 5~October 2026 (repository created 17:16~UTC). The
Atlas files \path{MongeCircles.lean} and \path{NapoleonTheorem.lean} first became public on
5~October 2026 between 21:37 and 21:40~UTC, when a batch of import pull requests (\#46--\#65) was
opened; this is after our work. These are dates of \emph{public} availability only. The Atlas
pull requests describe the files as migrated from an upstream repository,
\path{fairinternal/formal-math}, which is private, so we cannot see when the Atlas files were
written, and earlier private authorship cannot be excluded. Likewise, our own date of
3~October 2026 is a private date that readers cannot check; only the public dates are
verifiable. We therefore claim no priority over the Atlas proofs, only that ours was not derived
from them. The ShadowBench proof (12~June 2026) is earlier on any reading. Our searches
(GitHub code search for ``monge'' in Lean, Mathlib's \path{docs/1000.yaml}, which lists the
theorem with no declaration, Mathlib pull requests, the AFP, HOL~Light, Coq and Mizar) missed the
first of these; they were identified during revision. What remains distinguishable in our version
is narrow: the statement is in coordinates $K\times K$ over an arbitrary field~$K$, including
fields of positive characteristic, rather than over~$\RR$ or a real module, and
Section~\ref{sec:monge} analyses which hypotheses are needed and why (pairwise distinct radii only,
none on the centres, and the meaning of zero, negative or equal radii). We do not claim that our
hypotheses are weaker than those of the Atlas proof, whose general form also assumes only distinct
weights. We did not find a formalisation of the three-circle theorem in other systems, but given
the above we do not regard that search as reliable.

\paragraph{Simson.} The Wallace--Simson theorem is \emph{not} new to formal mathematics, and the
first version of ours was weaker than the existing ones in one respect. The Isabelle AFP entry \emph{The
Wallace--Simson Line Theorem in Isabelle/HOL} by Ramos, Hulak and de Queiroz~\cite{AFPSimson}
(session \lean{Simson}, 2026) uses complex coordinates, as we do, and proves both the theorem and
its geometric converse: if the three feet are collinear, then the point lies on the circumcircle.
Guilhot's Rocq development of oriented angles~\cite{GuilhotAngles,Guilhot} proves the theorem
\lean{droite\_Simson} in file \path{point_Simson.v} as an equivalence between collinearity of the
feet and concyclicity, that is, with its converse, under the additional hypothesis
\lean{\textasciitilde{} colineaire (vec C A) (vec C M)} (the point $M$ is not on line $CA$). Possessing
the converse does not make its admissible-configuration domain uniformly stronger than ours: our
forward theorem excludes no position of the point on the circle; it assumes only distinct
vertices and the circle condition. HOL~Light's \path{Complex/grobner_examples.ml}
\cite{HOLGrobner} checks an algebraic (Chou-style) encoding of Simson's theorem with the complex
Gr\"obner-basis procedure \lean{COMPLEX\_ARITH}; the geometric versions in the same file are
commented out as too slow. The main files of our Lean development prove only the forward direction; their
\decl{Simson.area2_perpFoot_eq_zero_iff} is a determinant-form statement between local
expressions, not the geometric converse proved in Isabelle and Rocq. Version~2 adds the geometric
converse in a Layer-15 file (\decl{Simson.Classical.cov_A3_b}, Section~\ref{sec:simson}), after
those formalisations; no priority is claimed for it. In Lean, the theorem is not in Mathlib and not in
\path{docs/1000.yaml}; a GitHub code search for ``simson'' in Lean found one planning file listing
``SimsonLine -- Define the Simson line'' with no proof~\cite{EGPlan}, and otherwise unrelated hits.
The original project search missed the Isabelle and Rocq entries; they were identified during
revision. As far as we know, ours is the first Lean proof of the forward direction, and nothing
more is claimed.

\paragraph{Napoleon.} Napoleon's theorem has prior formalisations, and our contribution is not
novelty. A Lean~4 repository by Imran Tanr\i{}kolu~\cite{Tanrikolu} (Lean \texttt{v4.15.0}, 64-line
file \path{napoleantheorem.lean}, no \lean{sorry}) proves in complex coordinates that the three
Napoleon points satisfy
$\mathrm{IsEquilateral}(z_1,z_2,z_3):\ (z_1-z_2)^2+(z_2-z_3)^2+(z_3-z_1)^2=0$. This is the
algebraic criterion for an equilateral triangle; the repository does not state or prove the
equality of distances, and treats one rotation only. Its core is the same kind of
\lean{linear\_combination} certificate as~\eqref{eq:nap}. The AFP entry \emph{Napoleon's
Theorem} by Ramos, Hulak and de Queiroz~\cite{AFPNapoleon} (2026) formalises the theorem in complex
coordinates, including the exterior-side condition and both orientations of the rotation.
The benchmark repository \path{epfl-lara/icml-26-lean-challenges}~\cite{ShadowBench}, file
\path{autformalization/problems/geometry/L2/geo_gen_L2_006/ShadowBench/Source/Main.lean} at commit
\texttt{7541fcec} (12~June 2026), contains complete Lean proofs of both the inner and the outer
Napoleon distance equalities, including side conditions (around lines 269 and 355); this is a
second earlier Lean proof, and it already states equality of distances for both triangles.
A further, independent Lean proof was published \emph{after} ours: the repository
\path{facebookresearch/atlas-lean}~\cite{AtlasLeanNapoleon}, file
\path{MathlibExt/Geometry/Euclidean/NapoleonTheorem.lean}, added by pull request~\#63 at commit
\texttt{27a38cfdcea9e80a11c49ff8e2ce17650e264552} (5~October 2026, 21:40~UTC; 400 lines, no
\lean{sorry}), proves the outer Napoleon theorem in \lean{EuclideanSpace ℝ (Fin 2)}
(\lean{napoleon\_theorem\_general} and \lean{napoleon\_theorem}), with the apexes given by distance
equalities and the exterior condition by signed areas, and shows that the common distance is
positive. Our Napoleon development was made in early October 2026 and made public at
\path{github.com/jpbald93/napoleon} on 5~October 2026 (repository created 17:16~UTC), a few hours
before the Atlas file became public, and independently of it. As for Monge, this is a statement
about public dates only: the Atlas file is migrated from the private repository
\path{fairinternal/formal-math}, so earlier private authorship cannot be excluded, and we claim
no priority over it.
Guilhot's Coq development of plane geometry with oriented angles also treats Napoleon's
theorem (file \path{point_napoleon.v})~\cite{GuilhotAngles,Guilhot}. In Mathlib, \path{docs/1000.yaml} lists Napoleon's theorem with no
declaration. Relative to the Tanr\i{}kolu repository, our version states equality of distances,
proves both Napoleon triangles and the orientation lemma; relative to the ShadowBench proof and the
AFP entry, it adds nothing of substance beyond its hypothesis analysis.

\paragraph{Lami.} \path{docs/1000.yaml} lists Lami's theorem with no declaration. GitHub code
search for ``Lami theorem'' and ``lami'' in Lean (51 and 68 hits) returned only false positives,
repository search found one Java and one HTML project, and a search of Mathlib issues and pull
requests returned nothing. Mathlib has the law of sines
(\mdecl{EuclideanGeometry.law_sin}, \mdecl{InnerProductGeometry.sin_angle_mul_norm_eq_sin_angle_mul_norm})
but no statement about equilibrium of forces.

\paragraph{Summary.} We found no earlier formal proof of Lami's theorem, in Lean or in other
systems; given that our searches missed existing Monge, Simson and Napoleon work, this is a report
of what we found, not a claim of priority. Monge's three-circle theorem has two earlier
independent Lean proofs, one predating ours by four months; our version differs only in its
arbitrary-field formulation and hypothesis analysis. The Desargues
development is, as far as we know, the first Lean proof of the projective-plane formulation; the
theorem is formalised in at least five other systems, including a coordinate proof in HOL~Light,
and an unmerged Mathlib pull request proves the affine (parallel) form. The Simson development is
a Lean proof of the forward direction (version~2 adds the converse) of a theorem already formalised, with its converse, in
Isabelle and in Rocq. Napoleon's theorem is formalised in Lean (twice), Isabelle and Rocq
already. The searches are not exhaustive, and the Monge, Simson and Napoleon cases show that they
missed existing work.

\section{Verification}\label{sec:verify}

\emph{The battery results below were produced on the build machine on 5~October 2026. The two
author-check confirmations have since been supplied and their placeholders resolved; no
placeholders remain in the paper. The version-2 commits are published at the hashes given in the
availability paragraph. The Layer-14 and Layer-15 checks of
version~2 are described in Section~\ref{sec:l1415}.}

\paragraph{What has already been checked.} Each project report records a full \lean{lake build}
with 0 errors and 0 warnings; on the build machine the cited commits build with 0 errors and 0
warnings in 1{,}805 jobs (Desargues), 3{,}013 (Monge), 8{,}711 (Simson), 8{,}709 (Napoleon) and
2{,}428 (Lami). The project reports also record an empty grep of the library sources for \lean{sorry},
\lean{admit}, \lean{native\_decide}, \lean{axiom}, \lean{\#eval}, \lean{run\_cmd},
\lean{set\_option}, \lean{macro}, \lean{elab}, \lean{syntax}, \lean{notation} and
\lean{import Lean}, and \lean{\#print axioms} output, from a file kept outside the library, showing
only \lean{propext}, \lean{Classical.choice} and \lean{Quot.sound} for 17 (Desargues), 3 (Monge),
4 (Simson), 5 (Napoleon) and 15 (Lami) declarations; the gate below covers all cited
declarations, including the bridge lemmas. On 5~October 2026 the library sources,
\path{gate.sh}, \path{tests/tamper.sh} and Lake files of the cited commits were compared by SHA-256
with the copies on the build machine and agree. The five repositories were also rebuilt from fresh
clones of the published commits on a second machine with a different architecture and operating
system (arm64, macOS): the builds are warning-free, \path{gate.sh} passes in all five at the same
declaration counts, and \path{tests/tamper.sh} rejects all six defects. Python scripts in each
\path{code/} directory check the statements numerically: for Desargues, 3{,}000 random perspective
configurations over $\QQ$ and about 600{,}000 configurations over $\F_5$ including degenerate ones;
for Monge, 2{,}000 random configurations over $\QQ$; for Simson, 2{,}000 random configurations with
an off-circle negative control; for Napoleon, 1{,}000 random triangles with two negative controls;
for Lami, 100{,}000 random equilibrium triples. None of these scripts is part of the formal proof.

\paragraph{House battery.} The same battery as for the earlier AFM papers was run on the
five repositories (reports of 5~October 2026, 20:42--21:27 UTC). The build machine's copies are not
git checkouts, so the reports record no commit; their Lean sources, \path{gate.sh},
\path{tests/tamper.sh} and Lake files are SHA-256-identical to those of the commits cited above, so
the results apply to those commits.
\begin{itemize}
\item \emph{Gate.} \path{gate.sh} in each repository rejects the forbidden tokens, builds the
library, generates its own \lean{\#print axioms} query for every declaration cited in this paper,
and fails on any non-standard axiom.
Result: PASS, standard axioms only, in all five: 24 declarations (Desargues), 7 (Monge),
12 (Simson), 21 (Napoleon), 17 (Lami).
\item \emph{Tamper tests.} \path{tests/tamper.sh} plants six defects in scratch copies (a
\lean{sorry}; a split \lean{axiom}; a fake \lean{\#eval} report; a macro; a deleted headline
hypothesis; a wrapped extra axiom) and requires each to be rejected.
Result: 6/6 rejected in each of the five repositories.
\item \emph{Conclusion and hypothesis mutations.} Small edits to conclusions and hypotheses of
each cited theorem, with the proof kept, must fail to compile; survivors are classified by hand.
Result: mutants rejected: Desargues 57/58 (2 by heartbeat timeout), Monge 16/16, Simson 29/29
(4 by heartbeat timeout), Napoleon 47/47, Lami 44/44. The one survivor, in
\decl{Desargues.ex_nonvacuous}, replaces $\ne0$ by $\ne1$ in the last conjunct; the value there
is $[a'b'c']=58$, which is neither $0$ nor~$1$, so the mutant is an equivalent (still true)
statement and its survival is expected. The Desargues run also covered a scratch file on the
build machine that is not in the repository (\path{dt1.lean}, a copy of the bracket identity); it
contributes one mutant (rejected by timeout), one statement and one ablation to the counts below.
Without it the Desargues figures are 56/57 mutants rejected (1 by timeout), 10 hypothesis sets
probed, 50 ablations NEEDED and 20/23 statements certified.
\item \emph{Contradiction probe.} Automation tries to derive \lean{False} from the hypotheses of
each theorem, with a control that must compile.
Result: hypothesis sets probed: 11 (Desargues), 2 (Monge), 10 (Simson), 4 (Napoleon),
14 (Lami), the remaining statements having no hypotheses; 0 found contradictory.
\item \emph{Hypothesis ablation.} Each Prop hypothesis is replaced by \lean{True} in a renamed copy,
which must then fail to compile. This shows that the proof uses the hypothesis, not that the
weakened statement is false; for that, see the remarks in Sections~\ref{sec:desargues}--\ref{sec:lami}.
Result: NEEDED 51 (Desargues), 5 (Monge), 24 (Simson), 7 (Napoleon), 53 (Lami); UNUSED 0 and
SKIP 0 in all five; no failed controls. In particular the new hypotheses $r_1\ne0$, $r_1\ne r_2$
of \decl{Monge.ext_eq_lineMap_homothety} and $\mathrm{area2}(a,b,c)\ne0$ of
\decl{Simson.concyc_eq_zero_iff_cospherical} are used by their proofs.
\item \emph{Unused-variable linter.} Every library file is compiled with
\lean{linter.unusedVariables} enabled, with a positive control.
Result: 0 warnings, in the library files and in the scratch files of all five repositories.
\item \emph{Satisfiability.} For each hypothesis-bearing statement, a kernel-checked
\lean{example} exhibits concrete witnesses for all hypotheses. The non-vacuity examples of
Sections~\ref{sec:desargues}--\ref{sec:napoleon} already cover the headline statements of
Desargues, Monge, Simson and Napoleon; Lami has none in the library.
Result: certified/total: Desargues 21/24, Monge 5/5, Simson 14/16, Napoleon 19/22,
Lami 14/17; every new bridge lemma is certified. Unresolved, and not counted as success:
\decl{Desargues.desargues_mk}, \decl{Desargues.desargues_converse_mk},
\decl{Desargues.br_eq_zero_of_isCollinear_mk}; \decl{Simson.simson},
\decl{Simson.simson_area}; \decl{Napoleon.apex_equilateral}, \decl{Napoleon.key},
\decl{Napoleon.napoleon_general}; \decl{Lami.lami_sin_ne_zero}, \decl{Lami.lami_nondegenerate},
\decl{Lami.lami_oriented_converse}.
\item \emph{Bridge inventory.} Every local \lean{def} is listed with its bridge to a Mathlib notion,
if any (Section~\ref{sec:correct}).
Result (heuristic: a notion counts as bridged if some lemma mentioning it is found): Monge 2/2
bridged (\decl{Monge.det3}, \decl{Monge.ext}); Simson 4/4 (\decl{Simson.footParam},
\decl{Simson.perpFoot}, \decl{Simson.area2}, \decl{Simson.concyc}); Napoleon 6/6
(\decl{Napoleon.centroid}, \decl{Napoleon.apex}, \decl{Napoleon.signedArea}, \decl{Napoleon.Ga},
\decl{Napoleon.Gb}, \decl{Napoleon.Gc}); Lami has no local definitions. For Desargues the scanner
reports 11 notions, 8 bridged (the bracket and the example points $O,A,B,C,P,Q,R$) and 3 not
(\lean{exA'}, \lean{exB'}, \lean{exC'}). That report is an artefact of the scanner: it searches
for the notion's name with a word-boundary pattern \lean{\textbackslash b}\emph{name}\lean{\textbackslash b},
which cannot match a name ending in a prime, so \lean{exA'}, \lean{exB'}, \lean{exC'} are never
found even though \decl{Desargues.exA'_isCollinear}, \decl{Desargues.exB'_isCollinear} and
\decl{Desargues.exC'_isCollinear} mention them. By inspection, all 11 Desargues notions have
bridge lemmas. The two Monge and Simson bridges carry the hypotheses stated in
Sections~\ref{sec:monge} and~\ref{sec:simson} ($r_1\ne0$; $\mathrm{area2}(a,b,c)\ne0$), and
Simson's \decl{Simson.area2} is bridged in one direction only at the cited commits (the other
direction is in the Layer-15 file of version~2).
\end{itemize}
Automated checks cannot replace comparing the formal statements with their informal counterparts.
The author has checked each entry of Table~\ref{tab:thms} by hand.

\subsection{Layers 14 and 15}\label{sec:l1415}

Version~2 adds two further checks, in two new files of each library, \path{L14.lean} and
\path{L15Classical.lean} (for example \path{Desargues/L14.lean}, \path{Desargues/L15Classical.lean}).
The library files of the cited commits are unchanged (SHA-256). With the new files the gates pass
with 34 declarations (Desargues), 15 (Monge), 21 (Simson), 45 (Napoleon) and 23 (Lami), standard
axioms only, and \path{tests/tamper.sh} rejects all six defects in each repository (arm64, macOS).

\paragraph{Layer 14: Mathlib-only statements and non-vacuity.} Each of the six numbered theorems has a
Mathlib-only formulation in \path{L14.lean}, proved from the library theorem through bridges; for
Napoleon it covers only non-degenerate configurations (see below). A mechanical check of the elaborated types, also with all implicit and instance arguments
shown, found no constant of the library's namespace. Each restatement has a witness theorem that
proves every hypothesis on an explicit configuration and that the configuration is non-degenerate:
\begin{itemize}
\item Theorem~\ref{thm:desdet}: \decl{Desargues.l14_desargues_det} (bracket as \mdecl{Matrix.det},
Mathlib's cross product); witness \decl{Desargues.l14w_desargues_det}: both triangles
non-degenerate, the three Desargues points computed and pairwise non-proportional.
\item Theorem~\ref{thm:desmk}: \decl{Desargues.l14_desargues}, quantified over points of
$\PP(K^3)$ rather than representatives, perspective as \mdecl{Projectivization.IsCollinear}; the
hypotheses ``representative $\ne0$'' disappear. Witness \decl{Desargues.l14w_desargues} over $\QQ$:
neither triangle collinear, $O$ distinct from the six vertices, $A\ne A'$, $B\ne B'$, $C\ne C'$, the
three Desargues points pairwise distinct.
\item Theorem~\ref{thm:monge}: \decl{Monge.l14_monge}, each centre of similitude given by its defining
property $r_1(O_2-S)=r_2(O_1-S)$, conclusion \mdecl{Collinear}, with exactly the paper's hypotheses
(and a homothety form \decl{Monge.l14_monge_homothety}, which needs nonzero radii for the ratio);
witness \decl{Monge.l14w_monge}: the example of Section~\ref{sec:monge}, centres not collinear,
the three centres of similitude pairwise distinct.
\item Theorem~\ref{thm:simson}: \decl{Simson.l14_simson}, feet as
\mdecl{EuclideanGeometry.orthogonalProjection}; witness \decl{Simson.l14w_simson}: $a=5$, $b=3+4i$,
$c=-4+3i$, $p=-5i$, triangle not collinear, $p$ not a vertex, the three feet computed and pairwise
distinct.
\item Theorem~\ref{thm:napoleon}: \decl{Napoleon.l14_napoleon_outer} and
\decl{Napoleon.l14_napoleon_inner}, apexes by equal sides and strict opposite (same) side, centres
by \mdecl{Finset.centroid}, no orientation hypothesis; witness \decl{Napoleon.l14w_napoleon}:
$A=0$, $B=2$, $C=2i$, the three centres computed and pairwise distinct.
\item Theorem~\ref{thm:lami}: \decl{Lami.l14_lami}, the library statement being already Mathlib-only
(and \decl{Lami.l14_lami_oriented}); witness \decl{Lami.l14w_lami}: forces $3$, $4i$, $-3-4i$,
pairwise distinct and nonzero, at a right angle, common ratio~$5$.
\end{itemize}
Layer~14 gives Mathlib-only formulations associated with all six numbered theorems. For Napoleon,
the strict-side formulation covers non-degenerate geometric configurations in either orientation; it
is not a restatement of the degenerate cases allowed by Theorem~\ref{thm:napoleon}, which remain
covered by the original rotation-based statements. The secondary
forms (\decl{Napoleon.l14_napoleon_inner}, \decl{Lami.l14_lami_oriented},
\decl{Monge.l14_monge_homothety}, \decl{Simson.l14_simson_feet}) have no separate witness, and the
converses of Table~\ref{tab:thms} are not restated in Layer~14 (the Desargues converse is covered
by Layer~15). The witnesses give concrete configurations for the Mathlib-only forms of
\decl{Desargues.desargues_mk}, \decl{Simson.simson} and \decl{Lami.lami_nondegenerate}, which the
satisfiability check above left unresolved; the battery figures above refer to the cited commits.

\paragraph{Layer 15: comparison with the classical statements.} Another AI agent collected
textbook statements of the five theorems from the literature, with citations (Coxeter;
Coxeter--Greitzer; Johnson; encyclopedic sources): Desargues and its converse (A1), Monge with its
companion form (A2), Wallace--Simson with its converse (A3), the outer and inner Napoleon
triangles (A4) and Lami (A5). An AI model given only those statements, which saw no code and not
this paper, wrote them in Lean with Mathlib notions, with its own definitions only where Mathlib has
none (centres of similitude and equilateral apexes by their defining properties): 9 theorems and 5
construction obligations, all with \lean{sorry}. Its transcript was audited: a transcript audit in the verification repository found no file read
beyond its own brief and outputs and no web access. It marked four
ambiguities and one extra hypothesis (positive radii, since Mathlib's spheres allow radius
$\le0$), and stated A3 in \lean{EuclideanSpace ℝ (Fin 2)}. A comparer then proved each of the 14
statements, copied verbatim, from the libraries at the cited commits without changing them, or
classified it: 9 covered by the library theorems (A1(a), A1(b), A2(a), A2(b), A3(a), A4(a), A4(b),
A5 and its obligation); 4 outside the paper's statements but proved (the obligations of A1 and the
two of A2, and the Simson converse A3(b)); none narrower than the library; 1 writer statement false
as written; none open. The writer's hypotheses $A\ne A'$, $B\ne B'$, $C\ne C'$ in A1(a) and distinct centres in A2(a)
were not needed; the latter agrees with Table~\ref{tab:nondeg}. New proofs, not in the main files:
\begin{itemize}
\item \decl{Simson.Classical.cov_A3_b}: the geometric converse of the Wallace--Simson theorem, through
the new bridge \decl{Simson.Classical.area2_eq_zero_of_collinear} (Section~\ref{sec:simson}).
\item \decl{Napoleon.Classical.cov_A4_a}: apexes by Mathlib's half-planes, and the outer Napoleon
triangle is affinely independent, by an identity for its doubled signed area.
\item \decl{Monge.Classical.cov_A2_b}: the companion form as a declaration; the paper had only the
instantiation $r\mapsto-r$ in prose.
\item \decl{Desargues.Classical.cov_A1_obligation}: two distinct lines of $\PP^2(K)$ meet in exactly
one point.
\item \decl{Monge.Classical.cov_A2_obligation_ext}, \decl{Monge.Classical.cov_A2_obligation_int}:
existence and uniqueness for two circles with positive radii and distinct centres, with unequal
radii additionally required for the external centre.
\item \decl{Napoleon.Classical.cov_A4_obligation_corrected}: on each side of a segment with $U\ne V$
there is exactly one equilateral apex.
\end{itemize}
Narrower cases: none. Writer mis-statement (one): \lean{classical\_A4\_obligation} omits $U\ne V$;
for $U=V$ the only apex is $U$ itself, on the ``line'', so no apex is strictly on either side. The
negation is proved (\decl{Napoleon.Classical.A4_obligation_false}) and the corrected statement
above. This is an error of the writer, not of the paper. The comparer also noted, as
Section~\ref{sec:napoleon} says, that the library's $\omega_o$ apexes are outward only for
counter-clockwise triangles, so the clockwise case of A4 uses \decl{Napoleon.napoleon_inner}. All
these proofs depend only on the three standard axioms. They are collected in
\path{L15Classical.lean} of each repository (namespace \lean{Desargues.Classical} and so on), whose
gate lists every one of them; the classical definitions and statements are reproduced there
verbatim.

\section{Limitations and further work}\label{sec:limits}

\begin{itemize}
\item \emph{Monge stops at a determinant.} The conclusion of \decl{Monge.monge} is $\det_3=0$ for
explicit formulas. The bridges \decl{Monge.det3_eq_zero_iff_collinear} and
\decl{Monge.ext_eq_lineMap_homothety} identify $\det_3=0$ with \mdecl{Collinear} and, for
$r_1\ne0$, $r_1\ne r_2$, the point \decl{Monge.ext} with the centre of a homothety sending $O_1$
to $O_2$, but the main files give no composed statement about circles: there are no circles in the
statement, and that the homothety maps one circle onto the other is not proved there. Version~2 adds
\decl{Monge.l14_monge}, with conclusion \mdecl{Collinear}, and proves in the Layer-15 file, for
circles with distinct centres and positive unequal radii given as Mathlib spheres, that the
homothety maps one circle onto the other at a unique external centre and, for three circles with
pairwise distinct radii, that the three external centres are collinear
(\decl{Monge.Classical.cov_A2_obligation_ext}, \decl{Monge.Classical.cov_A2_a}). All equal-radius cases are
excluded (one or three centres of similitude at infinity, or no unique centre for identical
circles; Section~\ref{sec:monge}), and the tangent-line formulation, which needs the external
tangents to exist, is not proved.
\item \emph{Converses.} In the main files the Desargues converse concludes collinearity in the dual
plane, and turning this into the existence of a point on three lines requires a second lemma; the
composition is proved in the Layer-15 file (\decl{Desargues.Classical.cov_A1_b}). The main files
prove only the forward direction of the Simson theorem, and
\decl{Simson.area2_perpFoot_eq_zero_iff} is only a determinant-level statement; the geometric
converse (proved in the Isabelle and Rocq developments) was not proved in the first version and is
proved in version~2, in the Layer-15 file (\decl{Simson.Classical.cov_A3_b}). The Lami converse is oriented only, as it must be.
\item \emph{Fields.} In the main files Simson and Napoleon are stated in $\CC$ only, not in a general Euclidean
plane (\lean{EuclideanSpace ℝ (Fin 2)} or an oriented two-dimensional affine space); the Layer-15
file proves the classical Simson statements (both directions) and both Napoleon statements in
\lean{EuclideanSpace ℝ (Fin 2)} (Section~\ref{sec:l1415}), and Lami's
oriented form needs a two-dimensional oriented space. The Simson certificate itself is a
polynomial identity; a version over an arbitrary ordered field, or over an arbitrary field with
non-isotropic side vectors, has not been attempted.
\item \emph{Orientation in Napoleon.} In the main files ``outward'' is expressed by the sign of the area of the
erected triangles, not by half-planes; the Layer-14 and Layer-15 statements use Mathlib's
half-planes (\mdecl{AffineSubspace.SOppSide}; Section~\ref{sec:l1415}).
\item \emph{Desargues in higher dimension.} Only the plane is treated; no statement in projective
space of dimension at least three, and no axiomatic version.
\item \emph{Lami and the law of sines.} The force-triangle derivation is not formalised.
\item \emph{Novelty.} Monge's three-circle theorem already has two complete Lean proofs, one of
them earlier than ours; Napoleon's theorem is formalised elsewhere, including in Lean; the
Wallace--Simson theorem and its converse are formalised in Isabelle and Rocq; Desargues's theorem
is formalised in several other systems (including by coordinates in HOL~Light), and its affine
parallel form is the subject of an open Mathlib pull request.
\end{itemize}

\section*{Acknowledgements and use of AI}
A large-language-model AI assistant was used to write the Lean proofs and the Python checking
scripts, to run the prior-art searches, and to draft this paper. The author directed the work and
reviewed the paper. The author checked each formal statement in Table~\ref{tab:thms} against its
informal counterpart by hand. The Layer-14 and Layer-15 files of version~2, the collection of the
classical statements, their Lean formulation and the comparisons were carried out by AI agents and
are not part of the author's hand check. The resulting
formal proofs and counterexample are checked by the Lean kernel; fidelity to the textbook sources
and the comparison classifications require separate semantic review.
Every proof is checked by the Lean kernel, and the author takes responsibility for all content.

\begin{thebibliography}{99}\sloppy
\bibitem{AFP} \emph{Archive of Formal Proofs}, \url{https://www.isa-afp.org/}, accessed 4 October
2026.
\bibitem{AFPProjective} A.~Bordg, Projective Geometry, \emph{Archive of Formal Proofs}, 2018,
\url{https://isa-afp.org/entries/Projective_Geometry.html}, accessed 4 October 2026.
\bibitem{AtlasLean} \emph{atlas-lean}, file
\path{MathlibExt/Geometry/Euclidean/MongeCircles.lean}, pull request \#62, commit
\texttt{39febc62dfe68f9651c72eedf7916a7a070141a8} (5 October 2026),
\url{https://github.com/facebookresearch/atlas-lean}, accessed 7 October 2026.
\bibitem{AtlasLeanNapoleon} \emph{atlas-lean}, file
\path{MathlibExt/Geometry/Euclidean/NapoleonTheorem.lean}, pull request \#63, commit
\texttt{27a38cfdcea9e80a11c49ff8e2ce17650e264552} (5 October 2026),
\url{https://github.com/facebookresearch/atlas-lean}, accessed 7 October 2026.
\bibitem{AFPSimson} A.~F.~Ramos, D.~B.~Hulak and R.~J.~G.~B.~de Queiroz, The Wallace--Simson Line
Theorem in Isabelle/HOL, \emph{Archive of Formal Proofs}, August 2026, session \texttt{Simson},
\url{https://isa-afp.org/entries/Simson.html}, accessed 4 October 2026.
\bibitem{Chou} S.-C.~Chou, \emph{Mechanical Geometry Theorem Proving}, D.~Reidel, Dordrecht, 1988.
\bibitem{Lean4} L.~de Moura and S.~Ullrich, The Lean~4 theorem prover and programming language, in
\emph{Automated Deduction -- CADE~28}, Lecture Notes in Computer Science 12699, Springer, 2021,
625--635.
\bibitem{EGPlan} \emph{EG}, file \path{Plan_files/Plan.lean},
\url{https://github.com/jjdishere/EG}, accessed 3 October 2026.
\bibitem{Guilhot} F.~Guilhot, Proofs with Coq of theorems in plane geometry using oriented angles,
INRIA research report, 2002, \url{https://inria.hal.science/inria-00072226v1}, accessed 4 October
2026.
\bibitem{GuilhotAngles} F.~Guilhot, \emph{Formalization of the oriented angles theory} (Coq
contribution \texttt{Angles}, 2002; files \path{point_Simson.v}, \path{point_napoleon.v}),
\url{https://github.com/rocq-archive/angles}, accessed 4 October 2026.
\bibitem{HOLDesargues} J.~Harrison, \emph{HOL Light}, file \path{100/desargues.ml} (Desargues's
theorem, \lean{DESARGUES\_DIRECT}), \url{https://github.com/jrh13/hol-light}, accessed 4 October
2026.
\bibitem{HOLGrobner} J.~Harrison, \emph{HOL Light}, file \path{Complex/grobner_examples.ml}
(``Simson's theorem (Chou, p7)''), \url{https://github.com/jrh13/hol-light}, accessed 4 October
2026.
\bibitem{MNS} N.~Magaud, J.~Narboux and P.~Schreck, A case study in formalizing projective
geometry in Coq: Desargues theorem, \emph{Computational Geometry} 45 (2012), no.~8, 406--424.
\bibitem{Mathlib} The mathlib Community, The Lean mathematical library, in \emph{Proceedings of the
9th ACM SIGPLAN International Conference on Certified Programs and Proofs (CPP 2020)}, ACM, 2020,
367--381. Repository: \url{https://github.com/leanprover-community/mathlib4}, accessed 6 October
2026.
\bibitem{PR37456} Mathlib pull request \#37456, \emph{feat(LinearAlgebra/AffineSpace): add
parallel form of Desargues's theorem}, opened 31 March 2026, open and unmerged at the time of
access, head commit \texttt{0900538cdc1e2951e2f6adb7298f5f8922612233},
\url{https://github.com/leanprover-community/mathlib4/pull/37456}, accessed 4 October 2026.
\bibitem{AFPNapoleon} A.~F.~Ramos, D.~B.~Hulak and R.~J.~G.~B.~de Queiroz, Napoleon's Theorem,
\emph{Archive of Formal Proofs}, 2026, \url{https://isa-afp.org/entries/Napoleon.html}, accessed
4 October 2026.
\bibitem{ShadowBench} \emph{icml-26-lean-challenges} (ShadowBench), files
\path{autformalization/problems/geometry/L2/geo_gen_L2_005/ShadowBench/Source/Main.lean} (Monge) and
\path{autformalization/problems/geometry/L2/geo_gen_L2_006/ShadowBench/Source/Main.lean} (Napoleon),
commit \texttt{7541fcecef97cc124f38355a2cd2ba0ec8485dd5} (12 June 2026),
\url{https://github.com/epfl-lara/icml-26-lean-challenges}, accessed 7 October 2026.
\bibitem{Tanrikolu} I.~Tanr\i{}kolu, \emph{Napoleon-s-Theorem-in-Lean4}, file
\path{napoleantheorem.lean}, commit \texttt{0e6eb11}, \url{https://github.com/imran-tanrikolu/Napoleon-s-Theorem-in-Lean4},
accessed 4 October 2026.
\bibitem{Wiedijk100} F.~Wiedijk, \emph{Formalizing 100 Theorems},
\url{https://www.cs.ru.nl/~freek/100/}, accessed 4 October 2026.
\end{thebibliography}

\end{document}
